\documentclass[reqno,a4paper,12pt]{amsart}
\usepackage[OT1]{fontenc}
\usepackage{amsmath,amssymb,amsthm}
\usepackage[textwidth=152mm,textheight=230mm,centering]{geometry}
\usepackage{microtype,xurl}
\usepackage[hidelinks]{hyperref}
\hypersetup{pdftitle={Schauder bases in unitary orbits: a modular proof}}
\newtheorem{theorem}{Theorem}[section]
\newtheorem{proposition}[theorem]{Proposition}
\newtheorem{lemma}[theorem]{Lemma}
\newtheorem{corollary}[theorem]{Corollary}
\theoremstyle{definition}
\newtheorem{definition}[theorem]{Definition}
\newtheorem{remark}[theorem]{Remark}
\newcommand{\R}{\mathbb R}
\newcommand{\C}{\mathbb C}
\newcommand{\Q}{\mathbb Q}
\newcommand{\Z}{\mathbb Z}
\newcommand{\N}{\mathbb N}
\newcommand{\T}{\mathbb T}
\newcommand{\ip}[2]{\langle #1,#2\rangle}
\newcommand{\norm}[1]{\left\lVert #1\right\rVert}
\newcommand{\spn}{\operatorname{span}}
\newcommand{\dist}{\operatorname{dist}}
\newcommand{\diag}{\operatorname{diag}}
\newcommand{\ran}{\operatorname{ran}}
\newcommand{\PB}{\operatorname{PB}}
\newcommand{\Id}{\operatorname{Id}}
\newcommand{\E}{\mathbb E}
\newcommand{\ml}[1]{\nolinkurl{#1}}
\title[Schauder bases in unitary orbits]{Schauder bases in unitary orbits\\A modular proof for formalization}
\date{November 18, 2026}
\begin{document}
\begin{abstract}
We prove the orbit-basis characterization of simple continuous spectrum.
The proof is organized around finite prefix inequalities, explicit
weighted products, finite-dimensional matrix estimates, and compatible
ordered extensions. Every construction beyond the enumerated Mathlib
foundations is proved in the document.
\end{abstract}
\maketitle
\raggedbottom
\begingroup\footnotesize
\tableofcontents
\endgroup
\clearpage

\section{The target and the foundational boundary}
\subsection{Statement}
Let $U=(U_t)_{t\in\R}$ be a strongly continuous one-parameter group of
unitary operators on a separable infinite-dimensional complex Hilbert
space $H$. A vector $v$ is \emph{cyclic} when
$\spn\{U_tv:t\in\R\}$ is dense. A \emph{common eigenvector} is a nonzero
$a$ such that $U_ta\in\C a$ for every $t$. The group has \emph{simple
continuous spectrum} when it has a cyclic vector and no common eigenvector.
A Schauder basis means a sequence having a unique norm-convergent ordered
scalar expansion for every vector.

\begin{theorem}\label{thm:main}
The group $U$ has simple continuous spectrum if and only if there are
$g\in H$ and distinct $t_n\in\R$ such that $(U_{t_n}g)_{n\ge1}$
is a Schauder basis of $H$.
\end{theorem}

\subsection{Conventions}
Inner products are conjugate-linear in the first argument and linear in
the second, as in Mathlib. All operator norms are Hilbert operator norms.
A finite matrix acts on Euclidean coefficient space, not on a space with
the coordinate supremum norm. Matrix inequalities are Loewner
inequalities: $A\le B$ means that $B-A$ is Hermitian and
$c^*(B-A)c\ge0$ for every $c$. Scalar comparisons of quadratic forms
refer to their real values. The same order convention applies to
bounded self-adjoint operators. Orthogonal direct sums always have the
sum-of-squares norm.

Finite ordered families are indexed by $0,\ldots,d-1$; an augmented
scalar family has its constant at index $0$ and its nonconstants at
$1,\ldots,d$. Infinite auxiliary sequences are indexed by $\N=\{0,1,\ldots\}$.
Shifting the index by one gives the indexing in Theorem~\ref{thm:main}.
We write $e_t(x)=\exp(2\pi itx)$. On $\T=\R/\Z$ and its finite powers
we use normalized Haar measure.

A finite Borel measure on $\R$ is called \emph{singleton-null} if every
singleton has measure zero. This is the meaning of atomlessness needed
in the construction, and corresponds to Mathlib's
\ml{NullSingletonClass}; no assertion about atoms in arbitrary measurable
spaces is required. Inequalities between general measurable weights
are understood almost everywhere unless explicitly stated pointwise.
The finite trigonometric weights constructed below obey their stated
pointwise bounds.

\subsection{What is imported from Mathlib}
Self-containment here is relative to the Mathlib foundations permitted
in the problem. We use:
\begin{enumerate}
\item Hilbert-space completeness, orthogonal projection, the
orthogonal-complement density criterion, Hilbert duality, and the
bounded inverse theorem for Banach spaces;
\item finite-dimensional self-adjoint diagonalization and the basic
algebra of finite matrices and continuous linear maps;
\item Lebesgue and Bochner integration, Fubini--Tonelli, dominated and
monotone convergence, and integration against a density;
\item the Gaussian Fourier formula, Fourier inversion for a continuous
$L^1$ function with $L^1$ Fourier transform, and uniqueness of
characteristic functions of finite positive measures on $\R$;
\item tightness of probability measures whose characteristic functions
converge to a function continuous at zero, compactness of a tight
family's closure, and the definition of weak convergence;
\item independence of distinct monoid characters, the analytic identity
principle, existence of an accumulation point in a positive-measure
set for a singleton-null measure on $\R$, and totality of the circle
characters in Haar $L^2$.
\end{enumerate}
The precise source revision and declarations are recorded in
Section~\ref{sec:mathlib}. Elementary real analysis, finite sums,
polynomial root bounds, and the topology of $\R$ are also used.
All specialized weighted-system, matrix-prefix, spectral-measure,
and limiting-basis arguments are supplied below. No FLT Lean code or
external weighted-basis theorem is an input.

\section{Representatives, weights, and ordered families}
\subsection{Finite characters}
\begin{definition}[Mixed characters]\label{def:characters}
For a finite coordinate set $I$, put $X_I=\R\times\T^I$.
For $(b,k)\in\Q\times\Z^I$ set
\[
 \chi_{b,k}(x,y)=e_b(x)\prod_{i\in I}\exp(2\pi i k_i y_i).
\]
A \emph{character polynomial} is a finite complex linear combination
of these functions. On a pure torus omit the real coordinate and $b$.
We retain these actual bounded continuous functions as representatives,
and only then take their classes in a specified $L^2$ space.
\end{definition}

\begin{lemma}[Character algebra and Haar moments]\label{lem:character-algebra}
Characters have modulus one, their products add frequencies, and their
conjugates negate frequencies. Character polynomials are closed under
addition, multiplication and conjugation, are bounded, and have finite
auxiliary Fourier support. On a pure torus,
\[
 \int\chi_k\,dy=\mathbf1_{\{k=0\}},\qquad
 \int\overline{\chi_k}\chi_l\,dy=\mathbf1_{\{k=l\}}.
\]
\end{lemma}
\begin{proof}
The exponential addition formula proves the algebraic assertions.
A finite sum $\sum a_j\chi_j$ is bounded by $\sum|a_j|$.
On one circle the integral over $[0,1]$ of $\exp(2\pi iny)$ is one
for $n=0$ and, by its antiderivative, zero for nonzero integer $n$.
Fubini gives the finite-product formula, including the empty product.
The second identity follows by conjugation and multiplication.
\end{proof}

\begin{lemma}[Weighted measures and their norms]\label{lem:weighted-measures}
Let $\mu$ be finite and $w$ measurable with $0<w\le1$ almost everywhere.
Then $w\mu$ is finite, has the same null sets as $\mu$, and is nonzero
if $\mu\ne0$. It is singleton-null if $\mu$ is. For bounded measurable
$f,g$ define
\[
 \ip{f}{g}_{w\mu}=\int\overline f g w\,d\mu,\qquad
 \|f\|_{w\mu}^2=\int |f|^2w\,d\mu.
\]
These are the inner products and norms of their $L^2(w\mu)$ classes.
If $w_j\to w$ in $L^1(\mu)$, the inner products of any fixed finite
list of bounded functions converge entrywise.
\end{lemma}
\begin{proof}
Finiteness and the displayed identities follow from integration against
a density. If $\int_Aw\,d\mu=0$, then
$\mu(A\cap\{w\ge1/n\})=0$ for every positive integer $n$.
Their union covers $A$ outside the null set $\{w=0\}$, proving equality
of null sets; all other measure assertions follow.
Finally
\[
 \left|\int\overline f g(w_j-w)\,d\mu\right|
 \le\|f\|_\infty\|g\|_\infty\|w_j-w\|_{L^1(\mu)}.
\]
\end{proof}

\begin{definition}[Prefix bound]\label{def:prefix}
For an ordered family $v=(v_0,\ldots,v_{d-1})$ in a Hilbert space and
$K\ge0$, write $\PB_K(v)$ if
\[
 \left\|\sum_{i<k}a_iv_i\right\|
 \le K\left\|\sum_{i<d}a_iv_i\right\|
 \quad(0\le k\le d,\ a\in\C^d).
\]
Its synthesis map is $X_va=\sum_i a_iv_i$, and its Gram matrix is
$G_v=(\ip{v_i}{v_j})_{ij}=X_v^*X_v$.
For a measure subscript we use the corresponding $L^2$ norm.
No optimal finite basis constant is needed.
\end{definition}

\begin{lemma}[Gram and prefix identities]\label{lem:gram-prefix}
For every $a$, $\|X_va\|^2=a^*G_va$. The family is independent exactly
when $G_v>0$. If $D_k$ is the Euclidean coordinate projection onto the
first $k$ entries, then
\[
 \PB_K(v)\quad\Longleftrightarrow\quad D_kG_vD_k\le K^2G_v
 \quad\hbox{for every }k.
\]
For a nonempty independent family a prefix bound must have $K\ge1$.
Multiplying the measure by a positive scalar leaves prefix bounds
unchanged and multiplies squared distances by that scalar.
\end{lemma}
\begin{proof}
Expand the squared norm as a double sum. Positivity is equivalent to
having zero kernel for synthesis. Apply the norm formula to $a,D_ka$
to get the matrix characterization. The full prefix is the identity
on a nonzero span, so $K\ge1$. The scaling claims follow by multiplying
each squared norm, and then taking the infimum for distance.
\end{proof}

\begin{lemma}[Three operations on prefix bounds]\label{lem:prefix-operations}
\begin{enumerate}
\item If $G>0$, $0<a\le b$, $aG\le L\le bG$, and $\PB_K(G)$ holds, then
$\PB_{\sqrt{b/a}K}(L)$ holds.
\item An ordered subfamily of a family with $\PB_K$ has $\PB_K$.
An orthogonal concatenation of families with bounds $K_i\le K$ has
bound $K$, provided $K\ge1$.
\item Let the first $r$ vectors of an independent family span $V$.
If all coordinate prefixes ending after that initial block have norm
at most $K$, the remaining vectors, orthogonally projected off $V$,
are independent and have $\PB_K$.
\end{enumerate}
\end{lemma}
\begin{proof}
For (1), $D_kLD_k\le bD_kGD_k\le bK^2G\le(b/a)K^2L$.
For a subfamily insert zero coefficients in the omitted coordinates.
For orthogonal concatenation, square the norm and sum the bounds in
the orthogonal components.
For (3), a zero projected tail combination lies in the initial span,
so independence makes its coefficients zero. For a tail combination
$x$ choose the initial vector $a=-P_Vx$. A prefix ending in the tail
acts on $a+x$ as $a+x_{\rm prefix}$. Its norm is at most
$K\|(I-P_V)x\|$. Projecting this output off $V$ gives the asserted
prefix inequality. Only prefixes ending after the initial block are used.
\end{proof}

\begin{lemma}[Concatenation across a small angle]\label{lem:angle}
Suppose nonempty independent families in $E,F$ have bounds $K_E,K_F$, and
$|\ip{x}{y}|\le\beta\|x\|\|y\|$ for $x\in E,y\in F$, with $0\le\beta<1$.
Their concatenation is independent and has bound
\[
 \sqrt{(1+\beta)/(1-\beta)}\max(K_E,K_F).
\]
\end{lemma}
\begin{proof}
The Cauchy--Schwarz bound on the cross term gives
\[
 (1-\beta)(\|x\|^2+\|y\|^2)\le\|x+y\|^2
 \le(1+\beta)(\|x\|^2+\|y\|^2).
\]
The lower bound proves independence. In the orthogonal direct-sum
metric the concatenation has bound $\max(K_E,K_F)$ by
Lemma~\ref{lem:prefix-operations}; the two inequalities transfer that
bound to the actual norm.
\end{proof}

\begin{lemma}[Distances under decreasing weights]\label{lem:distance-monotone}
Let $\mu$ be finite, let $f,A,A'$ consist of bounded measurable
functions, and let $0\le w'\le w\le1$ be measurable weights.
For a finite list $A$ put
$d_\nu(f,A)=\inf_a\|f-X_Aa\|_\nu$.
If $A'$ extends $A$ as an initial
segment, then $d_{w'\mu}(f,A')\le d_{w\mu}(f,A)$.
\end{lemma}
\begin{proof}
Pad each coefficient vector for $A$ by zeros. Its residual under $w'$
has norm at most the same residual under $w$. Take infima.
Thus this comparison does not identify two different $L^2$ spaces.
\end{proof}

\section{Matrix estimates with explicit metrics}
For a positive definite matrix $G$ we also write $\PB_K(G)$ for the
inequalities in Lemma~\ref{lem:gram-prefix}. Throughout this section
all matrix inverses are taken only for positive definite matrices.

\begin{lemma}[Inverse quadratic forms]\label{lem:inverse-variational}
If $G>0$, then
\[
 v^*G^{-1}v=\max_x\{2\operatorname{Re}(v^*x)-x^*Gx\}.
\]
Consequently inversion reverses positive matrix order. If $E$ is an
orthogonal projection, the inverse of $EGE$ on $\ran E$ is at most
$EG^{-1}E$ on that subspace. For a measurable positive matrix family
$G(s)\ge aI$, $a>0$, on a probability space, with $G$ integrable,
\[
 \left(\int G\,d\pi\right)^{-1}\le\int G^{-1}\,d\pi.
\]
\end{lemma}
\begin{proof}
Subtracting the expression in braces from the left side gives
$(x-G^{-1}v)^*G(x-G^{-1}v)$; its maximum is attained at $G^{-1}v$.
Increasing $G$ decreases each expression in braces. Restricting $x$
to $\ran E$ proves inverse compression. For averaging, each fixed $x$
satisfies
$2\operatorname{Re}(v^*x)-x^*G(s)x\le v^*G(s)^{-1}v$.
Integrate and take the maximum in $x$. The inverse family is measurable
by inverse continuity and bounded by $a^{-1}I$, so its integral exists.
Inverse continuity itself follows, near an invertible $G$, from
$B^{-1}-G^{-1}=B^{-1}(G-B)G^{-1}$ and a uniform positive lower bound.
\end{proof}

\begin{lemma}[Inverse compression along a dense exhaustion]
\label{lem:inverse-exhaustion}
Let $\mathcal H$ be a Hilbert space, $aI\le G\le bI$ with $a>0$,
and suppose a bounded inverse of $G$ is given. Let $V_n$ be increasing
finite-dimensional subspaces with dense union, all containing $v$.
If $G_n=P_{V_n}G|_{V_n}$, then
$\ip{v}{G_n^{-1}v}\uparrow\ip{v}{G^{-1}v}$.
\end{lemma}
\begin{proof}
Completing the square in the $G$ inner product gives the same
variational formula as Lemma~\ref{lem:inverse-variational}, now on
$\mathcal H$. Restricting to $V_n$ proves monotonicity and the upper
bound. Choose $x_n\in V_n$ tending to $G^{-1}v$: density and nesting
allow this by choosing successive approximants and retaining each until
the next one is available. The deficit of its variational value is
$\ip{x_n-G^{-1}v}{G(x_n-G^{-1}v)}\le b\|x_n-G^{-1}v\|^2\to0$.
The finite-dimensional maximum is at least this value.
\end{proof}

\begin{lemma}[Projection and distance from a Gram matrix]
\label{lem:gram-projection}
Let $X:\C^d\to\mathcal H$ be injective, $G=X^*X$, and $F=\ran X$.
Then $P_F=XG^{-1}X^*$. For $f\in\mathcal H$ and $g=X^*f$,
\[
 \dist(f,F)^2=\|f\|^2-g^*G^{-1}g.
\]
\end{lemma}
\begin{proof}
The candidate lies in $F$, and
$X^*(f-XG^{-1}X^*f)=0$, so its residual is orthogonal to $F$.
This characterizes the orthogonal projection. Pythagoras and
$\|XG^{-1}g\|^2=g^*G^{-1}g$ prove the distance formula.
\end{proof}

\begin{lemma}[Relative Gram stability]\label{lem:gram-stability}
If Hermitian $G_n\to G>0$ entrywise in a fixed positive finite dimension, then,
for every $0<\eta<1$, eventually
$(1-\eta)G\le G_n\le(1+\eta)G$. Thus a prefix bound $K$ for $G$
gives the bound $\sqrt{(1+\eta)/(1-\eta)}K$ for $G_n$.
If also $g_n\to g$ and $a_n\to a$, then
$a_n-g_n^*G_n^{-1}g_n\to a-g^*G^{-1}g$.
\end{lemma}
\begin{proof}
Choose $a_0>0$ with $G\ge a_0I$ by finite diagonalization.
In dimension $d$, $\|M\|\le d\max_{ij}|M_{ij}|$, as follows by
Cauchy--Schwarz in each row and summing the squared row bounds.
Thus $\|G_n-G\|\to0$. Eventually this norm is at most $\eta a_0$,
which bounds its quadratic forms between $-\eta G$ and $\eta G$.
Use Lemma~\ref{lem:prefix-operations} and the inverse identity from
Lemma~\ref{lem:inverse-variational} for the remaining assertions.
\end{proof}

\begin{lemma}[Operator norms]\label{lem:operator-norm}
For a linear map between finite-dimensional Hilbert spaces,
\[
 \|A^*\|=\|A\|,\qquad \|A^*A\|=\|AA^*\|=\|A\|^2.
\]
For a real matrix, its operator norms on real and complex Euclidean
space agree.
\end{lemma}
\begin{proof}
Write $\|A\|=\sup_{\|x\|,\|y\|\le1}|\ip{y}{Ax}|$ to prove the
adjoint identity. The equality
$\|Ax\|^2=\ip{x}{A^*Ax}$ and submultiplicativity give both inequalities
for $\|A^*A\|$; apply the result to $A^*$ for the other identity.
If $A$ is real and $z=x+iy$, then
$\|Az\|^2=\|Ax\|^2+\|Ay\|^2\le\|A\|_{\R}^2\|z\|^2$.
Restriction to real vectors proves the reverse inequality.
\end{proof}

\begin{lemma}[An oblique projection]\label{lem:oblique}
For a nonzero idempotent $P$ on a finite-dimensional Hilbert space,
\[
 \|P\|^2=1+\|P-P^*\|^2.
\]
If also $P\ne I$, then $\|I-P\|=\|P\|$.
\end{lemma}
\begin{proof}
Relative to $\ran P\oplus(\ran P)^\perp$,
$P=\left(\begin{smallmatrix}I&A\\0&0\end{smallmatrix}\right)$.
Its skew part has squared norm $\|A\|^2$ by squaring its adjoint
product. Also $PP^*=\diag(I+AA^*,0)$.
Diagonalizing $AA^*\ge0$ on the nonzero range gives
$\|PP^*\|=1+\|AA^*\|$, proving the formula by
Lemma~\ref{lem:operator-norm}. Apply the same formula to $I-P$
when it is nonzero; its skew part is the negative of that of $P$.
\end{proof}

\begin{lemma}[A coordinate projection through its two Gram matrices]
\label{lem:dual-gram}
Suppose $P=RS^*$ is a coordinate projection, where the columns of $R$
are its selected range vectors and those of $S$ their biorthogonal
dual vectors. If $C=R^*R>0$ and $D=S^*S$, then
$\|P\|^2=\|C^{1/2}DC^{1/2}\|$.
\end{lemma}
\begin{proof}
$V=RC^{-1/2}$ is an isometry onto $\ran R$, and
$PP^*=V(C^{1/2}DC^{1/2})V^*$.
Conjugation by this isometry preserves the norm on its range, and the
operator is zero on the orthogonal complement. Use
Lemma~\ref{lem:operator-norm}.
\end{proof}

\section{Explicit weighted products and information operators}
\begin{definition}[A constant-marginal product]\label{def:product}
Let $(S,\rho)$ be a probability space, let $Y$ be a finite torus, and
let $W:S\times Y\to(0,1]$ be measurable with
$\int_YW(s,y)\,dy=M>0$ for every $s$.
Set $d\sigma=M^{-1}W\,d\rho\,dy$ and, for bounded measurable $u$, set
\[
 \iota h(s,y)=h(s),\qquad
 \mathcal Eu(s)=M^{-1}\int_Yu(s,y)W(s,y)\,dy,\qquad
 u^\circ=u-\iota\mathcal Eu.
\]
All functions used for this operation below are bounded.
\end{definition}

\begin{lemma}[Marginals and centering]\label{lem:centering}
$\sigma$ is a probability, $\iota:L^2(\rho)\to L^2(\sigma)$ is an
isometry, $\mathcal Eu\in L^2(\rho)$, and $u^\circ$ is orthogonal to
$\iota L^2(\rho)$. In particular
$\|u\|_\sigma^2=\|\mathcal Eu\|_\rho^2+\|u^\circ\|_\sigma^2$.
\end{lemma}
\begin{proof}
Integrate $W/M$ first in $y$ to get mass one and the norm identity
for $\iota$. Fiberwise Cauchy--Schwarz gives
$|\mathcal Eu(s)|^2\le M^{-1}\int |u(s,y)|^2W(s,y)\,dy$,
so the mean is square integrable. Its subtraction makes the weighted
fiber integral zero. Fubini gives
$\ip{\iota h}{u^\circ}=0$ first for bounded $h$ and then for all $L^2$
functions by truncation and Cauchy--Schwarz. Pythagoras proves the last
identity. This also shows independence of the chosen a.e. representatives.
\end{proof}

\begin{definition}[Mean and centered Gram data]\label{def:information}
For a finite bounded family $(u_k)$ put
$Xc=\sum c_ku_k$, $Tc=\sum c_k\mathcal Eu_k$, and
$X_0c=\sum c_ku_k^\circ$. Then $X=\iota T+X_0$.
When $C=X_0^*X_0>0$, define $J=TC^{-1}T^*$ on $L^2(\rho)$.
This is a positive finite-rank operator with range in the span of
the means.
\end{definition}

\begin{lemma}[The information identity]\label{lem:information}
In Definition~\ref{def:information}, the full Gram matrix is
$G=C+T^*T>0$. If $F=\ran X$, then
\[
 \iota^*P_F\iota=T(C+T^*T)^{-1}T^*=\Psi(J),\qquad
 \Psi(J)=J(I+J)^{-1}.
\]
Moreover, for positive finite-rank $J,J'$,
$\|\Psi(J)-\Psi(J')\|\le\|J-J'\|$.
If $\|J\|\le a$, then $C\le G\le(1+a)C$.
\end{lemma}
\begin{proof}
Lemma~\ref{lem:centering} makes the two summands of $X$ orthogonal,
so $G=C+T^*T$ and $X^*\iota=T^*$. Apply
Lemma~\ref{lem:gram-projection}. Put $A=TC^{-1/2}$. The identity
$(I+AA^*)A=A(I+A^*A)$, followed by multiplication by inverses, gives
$A(I+A^*A)^{-1}A^*=AA^*(I+AA^*)^{-1}$.
Here the inverses on the base space are legitimate: restrict to the
finite span of the means, diagonalize the positive operator, and use
the identity on its orthogonal complement. Their norms are at most one.
The same finite-dimensional reduction, now on the sum of the ranges
of $J,J'$, justifies
\[
 \Psi(J)-\Psi(J')=(I+J)^{-1}(J-J')(I+J')^{-1}.
\]
Finally $\|TC^{-1/2}\|^2=\|J\|$ by
Lemma~\ref{lem:operator-norm}, applied on the finite span of its range.
Thus $T^*T\le aC$.
\end{proof}

\begin{lemma}[Independent auxiliary blocks]\label{lem:independent-blocks}
For finitely many weights $W_i(s,y_i)$ as in Definition~\ref{def:product},
use the density $\prod_i W_i/M_i$ over $\rho$ and the product torus.
Each block keeps its original mean and centered Gram matrix.
Centered vectors from different blocks are orthogonal. If their
centered Gram matrices are positive, the concatenated data satisfy
\[
 C=\diag(C_i),\qquad J=\sum_i T_iC_i^{-1}T_i^*.
\]
The concatenated centered prefix bound is the maximum of the
individual bounds.
\end{lemma}
\begin{proof}
Integrating all unused variables gives one. For two centered factors
in distinct variables, the fiber integral factors into their two
zero means. This proves orthogonality and the diagonal formula by
Fubini. Inverting that finite diagonal matrix proves the identity
for $J$. Apply Lemma~\ref{lem:prefix-operations} to the orthogonal
centered spans.
\end{proof}

\begin{lemma}[Joining small mean blocks]\label{lem:join}
Let $V_0,H_1,\ldots,H_L$ be finite-dimensional orthogonal Hilbert
spaces. Give each $H_i$ an ordered basis with bound $\sqrt Q$.
Let $T_i:H_i\to V_0$ satisfy
$\|\sum_iT_iT_i^*\|\le Z$ and $\|T_iT_i^*\|\le\eta$, where $Q>1+Z$.
Replace each vector $x$ of block $i$ by $x+T_ix$.
Every prefix containing all of $V_0$ has squared norm at most
$Q+Q^2\eta/(Q-1-Z)$.
\end{lemma}
\begin{proof}
For a cut inside block $i$, let $P$ be that block's prefix.
Completed $H_j$ components split off as identities. On the others the
operator is
$\left(\begin{smallmatrix}A&-T_i(I-P)\\0&P\end{smallmatrix}\right)$,
where $A$ sends $(v_0,(x_j)_{j>i})$ to $v_0-\sum_{j>i}T_jx_j$.
Hence $AA^*=I+\sum_{j>i}T_jT_j^*$ and $\|A\|^2\le1+Z$.
Also $\|P\|^2\le Q$ and $\|T_i(I-P)\|^2\le Q\eta$:
use Lemma~\ref{lem:oblique} for $P\ne0,I$ and check $0,I$ directly.
For $h=(Q-1-Z)/(1+Z)>0$, the scalar Cauchy--Schwarz inequality gives
$\|a-b\|^2\le(1+h)\|a\|^2+(1+h^{-1})\|b\|^2$.
It bounds the squared image of $(x,y)$ by
$Q\|x\|^2+(Q+Q^2\eta/(Q-1-Z))\|y\|^2$.
\end{proof}

\section{Scalar data and the sharp row estimate}
\subsection{Scalar data}
\begin{definition}[Measurable scalar data]\label{def:scalar}
Scalar data consist of a finite torus, a measurable even weight
$c_0\le w\le1$ pointwise for some $c_0>0$, and an ordered nonempty
list $\chi_1,\ldots,\chi_d$ of distinct nonconstant characters.
Evenness means $w(-y)=w(y)$. Write
\[
 \begin{gathered}
 m=\int w,\qquad\epsilon=1-m,\\
 B=\begin{pmatrix}m&b^*\\b&G\end{pmatrix},\qquad
 b_k=\int\overline{\chi_k}w,\\
 A=G-bb^*/m,\qquad C=A/m,\qquad a=b/m.
 \end{gathered}
\]
Here $B$ is the Gram matrix of $(1,\chi_1,\ldots,\chi_d)$ in
unnormalized $L^2(w)$; $C$ is the centered Gram matrix under $w/m$.
The gain is $Z=a^*C^{-1}a$.
Polynomial scalar data additionally require $w$ to be a character polynomial.
\end{definition}

\begin{lemma}[Scalar identities]\label{lem:scalar-identities}
For scalar data, $0<B\le I$, all its diagonal entries are $m$, and
$B$ is real symmetric. The matrices $A,C$ are positive definite.
For a polynomial with coefficient column $c$ the mean under $w/m$ is
$a^*c$ and its variance is $c^*Cc$. Moreover
\[
 mB^{-1}=
 \begin{pmatrix}1+Z&-a^*C^{-1}\\-C^{-1}a&C^{-1}\end{pmatrix},
 \qquad m(B^{-1})_{00}=1+Z.
\]
If $Z>0$, the least variance among polynomials of mean one is $1/Z$,
attained at $c=C^{-1}a/Z$. The squared norm of the projection of $1$
onto the nonconstant span under $w/m$ is $Z/(1+Z)$.
The squared norm of the first augmented coordinate prefix is $1+Z$.
\end{lemma}
\begin{proof}
Haar orthogonality and $c_0\le w\le1$ give $c_0I\le B\le I$.
Unimodularity gives its diagonal. Inverting all torus coordinates
conjugates every matrix entry while leaving its integral unchanged,
so the entries are real. Completing the first coordinate's square
in the quadratic form of $B$ proves $A>0$.
Expansion of the variance gives $C$. Multiplication verifies the
displayed inverse formula. Cauchy--Schwarz in the $C$ metric gives
$|a^*c|^2\le Zc^*Cc$, with equality at the displayed minimizer.
Since $(C+aa^*)C^{-1}a=(1+Z)a$, the projection formula of
Lemma~\ref{lem:gram-projection} gives $Z/(1+Z)$.
For the first augmented prefix use
Lemma~\ref{lem:dual-gram} with the one-by-one matrices $m$ and
$(B^{-1})_{00}$.
\end{proof}

\subsection{Finite-dimensional estimates}
\begin{lemma}[A real skew arrowhead]\label{lem:arrowhead}
For a real skew-symmetric $n$-by-$n$ matrix $S$, $n\ge1$, a real
column $b$, and $s\ge0$,
\[
 \left\|\begin{pmatrix}S&sb\\-sb^T&0\end{pmatrix}\right\|^2
 =\lambda_{\max}(-S^2+s^2bb^T).
\]
If $0\le r\le s$, the left-hand squared norm increases from $r$ to $s$
by at most $(s^2-r^2)\|b\|^2$.
\end{lemma}
\begin{proof}
Denote the skew block matrix by $T_s$. Real and complex norms agree by
Lemma~\ref{lem:operator-norm}. The Hermitian matrix $iT_s$ has spectrum
symmetric about zero, because complex conjugation sends it to $-iT_s$.
For $q>\|S\|^2$, put $h=\sqrt q$. Thus $\|T_s\|\le h$ is equivalent
to $hI-iT_s\ge0$. Its upper-left block $hI-iS$ is positive definite.
Completing the block square, or taking its Schur complement, gives
\[
 h-s^2b^T(hI-iS)^{-1}b\ge0.
\]
Since $(hI-iS)^{-1}=(hI+iS)(qI+S^2)^{-1}$, and
$S(qI+S^2)^{-1}$ is real skew, this is
$s^2b^T(qI+S^2)^{-1}b\le1$.
Conjugating by the square root of $qI+S^2>0$ makes this equivalent to
$qI+S^2-s^2bb^T\ge0$. This last step is the elementary fact
$I-zz^*\ge0$ exactly when $\|z\|^2\le1$.
Both squared norms being compared are at least $\|S\|^2$:
restrict $T_s$ to vectors $(x,0)$, and use $-S^2+s^2bb^T\ge-S^2$.
Their identical upper bounds for every $q>\|S\|^2$ prove the equality.
Finally the largest Rayleigh quotient increases by at most
$(s^2-r^2)\|b\|^2$ on adding that rank-one matrix.
\end{proof}

\begin{lemma}[A two-Gram interpolation bound]\label{lem:row-matrix}
Let $B$ be real positive definite with $B\le I$, put $H=B^{-1}$,
and suppose $B_{jj}=m$. Let $P=\{0,\ldots,j-1\}$ and
$F_{j-1}=\|B^{1/2}E_PB^{-1/2}\|^2$, except that $F_{-1}=1$.
Put $F_j=\|B^{1/2}E_{P\cup\{j\}}B^{-1/2}\|^2$.
For $0\le t\le1$ define
\[
 C_t=\begin{pmatrix}B_{PP}&\sqrt t\,B_{Pj}\\
 \sqrt t\,B_{jP}&1-t+tm\end{pmatrix},\quad
 D_t=\begin{pmatrix}H_{PP}&\sqrt t\,H_{Pj}\\
 \sqrt t\,H_{jP}&1-t+tH_{jj}\end{pmatrix}.
\]
Then
\[
 \|C_t^{1/2}D_tC_t^{1/2}\|
 \le\max(F_{j-1},F_j)+t(1-t)(1-m)(H_{jj}-1).
\]
\end{lemma}
\begin{proof}
Set $x=B^{1/2}e_j$, $y=B^{-1/2}e_j$,
$P_0=B^{1/2}E_PB^{-1/2}$ and $P_1=P_0+xy^T$.
Then $y^Tx=1$, $P_0x=0$, and $y^TP_0=0$.
With $s=\sqrt{t(1-t)}$, direct multiplication shows that
\[
 R_t=\begin{pmatrix}P_0+txy^T&sx\\sy^T&1-t\end{pmatrix}
\]
is a nonzero projection. Its range columns are
$(B^{1/2}e_k,0)$ for $k<j$ and $(\sqrt t\,x,\sqrt{1-t})$;
its dual columns replace $B^{1/2},x$ by $B^{-1/2},y$.
They are biorthogonal and independent even at $t=0,1$, and their
Gram matrices are $C_t,D_t$. Lemma~\ref{lem:dual-gram} identifies
the left side with $\|R_t\|^2$.

Let $S_i=P_i-P_i^T$, $A_0=S_1-S_0=xy^T-yx^T$ and $v=x-y$.
Put $a_0^2=mH_{jj}-1\ge0$, by Cauchy--Schwarz. Then
\[
 R_t-R_t^T=
 \begin{pmatrix}(1-t)S_0+tS_1&sv\\-sv^T&0\end{pmatrix},
 \quad
 \|v\|^2-a_0^2=(1-m)(H_{jj}-1)\ge0.
\]
If $a_0>0$, multiplication on $\spn(x,y)$ gives $A_0^2=-a_0^2I$.
Thus $u=-A_0v/(a_0\|v\|)$ is a unit vector and
$A_0u=a_0v/\|v\|$.
Conjugate $\diag(S_0,S_1)$ by
\[
 Q=\begin{pmatrix}\sqrt{1-t}\,I&\sqrt t\,I\\
 -\sqrt t\,I&\sqrt{1-t}\,I\end{pmatrix}.
\]
It is orthogonal. Compress its second block to $\R u$.
The resulting skew matrix has upper-left block $(1-t)S_0+tS_1$,
last column $sa_0v/\|v\|$, and norm at most
$\max(\|S_0\|,\|S_1\|)$.
By Lemma~\ref{lem:arrowhead}, enlarging that column to $sv$ costs
at most $s^2(\|v\|^2-a_0^2)$ in squared norm.
If $a_0=0$, equality in Cauchy--Schwarz gives $A_0=0$;
the same bound follows from that lemma starting with last column zero.
Now use Lemma~\ref{lem:oblique} for $R_t,P_1$, and for $P_0$ when
it is nonzero. When $P_0=0$, the sentinel value
$F_{-1}=1=1+\|S_0\|^2$ is a valid upper bound.
\end{proof}

\begin{lemma}[One finite-rank projection and one orthogonal projection]
\label{lem:two-projections}
Let $\Pi,D$ be orthogonal projections on a Hilbert space, with $\Pi$
of finite rank. The space decomposes orthogonally into two-dimensional
invariant blocks on which
\[
 \Pi=\begin{pmatrix}1&0\\0&0\end{pmatrix},\quad
 D=\begin{pmatrix}t&\sqrt{t(1-t)}\\
 \sqrt{t(1-t)}&1-t\end{pmatrix}\quad(0<t<1),
\]
and common eigenspaces with eigenvalues zero or one.
Only the common eigenspaces can be infinite dimensional.
\end{lemma}
\begin{proof}
The finite-dimensional space $\ran\Pi+D(\ran\Pi)$ is invariant under
both projections; self-adjointness also makes its orthogonal complement
invariant, and $\Pi=0$ there. Diagonalize $\Pi D\Pi$ on $\ran\Pi$.
For a unit eigenvector $v$ with eigenvalue $0<t<1$, let
$z=(Dv-tv)/\sqrt{t(1-t)}$. Then
$\Pi z=0$, $\|z\|=1$, and
$Dv=tv+\sqrt{t(1-t)}z$.
Applying $D$ again gives
$Dz=\sqrt{t(1-t)}v+(1-t)z$.
For distinct orthonormal eigenvectors $v_i,v_j$,
$\ip{Dv_i}{Dv_j}=\ip{v_i}{Dv_j}=t_j\delta_{ij}$,
which proves orthogonality of the resulting blocks.
Eigenvalues zero and one give respectively $Dv=0$ and $Dv=v$.
On the remaining invariant orthogonal complement $\Pi=0$; split it
as $\ran D\oplus\ker D$ to finish.
\end{proof}

\subsection{Activation on a rare set}
\begin{lemma}[A symmetric rare set]\label{lem:rare-set}
For $0<p<1$ the circle has a measurable symmetric set $D$ of Haar
measure $p$.
\end{lemma}
\begin{proof}
Take the image of $[-p/2,p/2]$ in $\R/\Z$. It is symmetric.
Identifying the circle with $[-1/2,1/2)$ computes its measure as its
interval length; endpoints are null.
\end{proof}

\begin{lemma}[Sharp activated row prefixes]\label{lem:row-split}
For scalar data and a symmetric $D\subset\T$ of measure $p$, put
$W(x,y)=1-\mathbf1_D(y)(1-w(x))$.
Let $Y$ be any finite list of distinct circle characters starting with
$1$. Order the characters $\chi_j\psi$ by $j=0,\ldots,d$ and then by
$\psi\in Y$, with $\chi_0=1$.
If $F_j$ is the squared norm of the old augmented prefix through
$\chi_j$ and $F_{-1}=1$, a new prefix ending in row $j$ has squared
norm at most
\[
 \max(F_{j-1},F_j)+\frac{\epsilon}{4}\bigl((B^{-1})_{jj}-1\bigr).
\]
The corresponding centered prefix has the same upper bound.
\end{lemma}
\begin{proof}
Use the Hilbert space of $(d+1)$-tuples of Haar $L^2(\T)$ functions.
Its activated norm has Gram operator
\[
 \mathcal G=I+(B-I)D,\qquad
 \mathcal G^{-1}=I+(B^{-1}-I)D,
\]
where $D$ multiplies every component by $\mathbf1_D$.
These formulas follow pointwise inside and outside $D$; in particular
the norm is equivalent to the ordinary sum-of-squares norm.
The synthesis $(F_j)_j\mapsto\sum_j\chi_j(x)F_j(y)$ identifies this
norm with the $W$-weighted norm by integration first in $x$.
On this full tuple space the prefix keeps all preceding rows, applies
the finite Fourier projection $\Pi$ onto the selected characters in
row $j$, and kills later rows. Its restriction is the prefix in question.

Apply Lemma~\ref{lem:two-projections} to $\Pi$ and multiplication by
$\mathbf1_D$ on $L^2(\T)$. The resulting decomposition is orthogonal
also for $\mathcal G$ and is invariant under the prefix.
Common eigenspaces give either an ordinary coordinate projection,
or an old $B$-metric prefix, so have squared bounds
$1,F_{j-1},F_j$.
On a two-dimensional block choose a unit vector in $\ran D$ and one
in $\ker D$, so the selected vector in $\ran\Pi$ has coordinates
$(\sqrt t,\sqrt{1-t})$. For each completed row the $\ker D$ coordinate
is an orthogonal identity summand, also in the dual Gram matrix.
Discard these identity summands. The remaining range and dual columns
have exactly the matrices $C_t,D_t$ in Lemma~\ref{lem:row-matrix}:
completed $\ran D$ coordinates give $B_{PP},(B^{-1})_{PP}$, the
selected row gives the factors $\sqrt t$ and diagonal $1-t+tm$.
Lemmas~\ref{lem:dual-gram} and \ref{lem:row-matrix} give the bound,
since $t(1-t)\le1/4$ and $F_j\ge1$.
Summing the squared bounds over the orthogonal blocks controls the
full tuple operator. Restrict to the finite product family, then apply
the quotient argument in Lemma~\ref{lem:prefix-operations} to its
initial constant to obtain the centered bound.
\end{proof}

\section{Constructing scalar data with arbitrarily small loss}
\subsection{Activation, truncation, and copying}
\begin{lemma}[Activated inverse bounds and gain limit]\label{lem:rare-data}
Use scalar data with $0<\epsilon<1$, $Z>0$ and activation set of
measure $0<p<1$. Put $M=1-p\epsilon$ and $H=B^{-1}$.
For a finite product family as in Lemma~\ref{lem:row-split}, let
$A_W$ be its unnormalized centered Gram matrix. A nonconstant product
in old row $j$ satisfies
\[
 (A_W^{-1})_{kk}\le1+p(H_{jj}-1),\qquad
 \|A_W^{-1}\|\le\|B^{-1}\|.
\]
For $Y_N=(1,e_1,e_{-1},\ldots,e_N,e_{-N})$ the product gain $z_N$
increases to the positive number
\[
 z_p=M[1+p(H_{00}-1)]-1
 =\frac p{1-\epsilon}
   \{Z+\epsilon^2-p\epsilon(Z+\epsilon)\}.
 \tag{A}\label{eq:rare-gain}
\]
For any $0<\alpha<1$ one can therefore choose a finite $Y_N$ with
$(1-\alpha)z_p\le z_N\le z_p$.
\end{lemma}
\begin{proof}
The finite augmented Gram matrix is a compression of $\mathcal G$
from Lemma~\ref{lem:row-split}. Its centered inverse is the
nonconstant principal block of its augmented inverse, by
Lemma~\ref{lem:scalar-identities}. Inverse compression bounds it by
the corresponding block of $\mathcal G^{-1}$.
For a coordinate character $\psi$, $\int_D|\psi|^2=p$, giving the
diagonal bound. Since $B\le I$, $\|B^{-1}\|\ge1$; the pointwise
inverse formula gives the norm bound.

The spaces consisting of tuples with each entry in $\spn Y_N$ are
nested and dense: use the Mathlib totality of the circle characters
and the finite number of tuple entries. The activated norm is
equivalent to the original one. Apply Lemma~\ref{lem:inverse-exhaustion}
to the unit coordinate vector corresponding to the constant in row
zero. Its full inverse quadratic form is $1+p(H_{00}-1)$.
Lemma~\ref{lem:scalar-identities} converts this into the gain limit.
Substitute $H_{00}=(1+Z)/(1-\epsilon)$ to obtain (A).
Its braces equal $Z(1-p\epsilon)+\epsilon^2(1-p)>0$.
Monotone convergence to this positive limit permits the finite choice.
\end{proof}

\begin{lemma}[Copies of scalar data]\label{lem:scalar-copies}
Take $L\ge1$ independent copies of scalar data with mass $M$,
normalized mean column $a$, centered matrix $C$, and gain $z$.
Use the product weight and concatenate the nonconstant lists.
The new data have
\[
 m'=M^L,\quad C'=\diag(C,\ldots,C),\quad
 a'=(a,\ldots,a),\quad Z'=Lz.
\]
Their centered bound is the old centered bound. If it is $\sqrt Q$
and $Q>1+Z'$, every augmented prefix has squared norm at most
$Q+Q^2z/(Q-1-Z')$.
\end{lemma}
\begin{proof}
All copied characters are nonconstant, and disjoint variable blocks
make their full frequencies distinct. The product weight is even and
bounded below by a positive constant. Integrating the unused
probability factors gives the displayed formulas, by
Lemma~\ref{lem:independent-blocks}.
For the last bound, regard the normalized constant functions as
$V_0$ and each centered span as $H_i$. The mean map on $H_i$ has
squared norm $z$ by Cauchy--Schwarz in its centered metric.
The sum of these rank-one mean operators has norm $Lz=Z'$.
Apply Lemma~\ref{lem:join}. The first prefix containing just the
constant is included in that lemma.
\end{proof}

\subsection{The numerical induction}
\begin{lemma}[Initial scalar data]\label{lem:scalar-initial}
For $0<\delta\le1$ put $a=\delta/400$ and
$w_0(y)=1-a+a\cos(2\pi y)$, with the single nonconstant character
$e^{2\pi iy}$. Then
\[
 \epsilon_0=a,\quad B_0=\begin{pmatrix}1-a&a/2\\a/2&1-a\end{pmatrix},
 \quad Z_0=\frac{a^2}{4-8a+3a^2},\quad
 a^2/4\le Z_0\le a^2/3.
\]
The centered bound is one,
$B_0\ge a^2I$, and $(B_0^{-1})_{11}\le1+2a$.
\end{lemma}
\begin{proof}
$1-2a\le w_0\le1$, and elementary Haar moments give the matrix.
Its eigenvalues are $1-a\pm a/2$, so the lower bound follows from
$a\le1/400$. Inverting the two-by-two matrix gives the gain and
\[
 (B_0^{-1})_{11}-1
 =\frac{a(1-3a/4)}{1-2a+3a^2/4}\le2a.
\]
The gain denominator lies between three and four.
A one-vector centered family has prefix bound one.
\end{proof}

\begin{lemma}[Choosing one quantitative step]\label{lem:scalar-numerics}
Suppose $0<\epsilon\le a\le1/400$ and
$3a^2/16\le Z\le Z_0$, and define $z_p$ by (A).
For any $C_0\ge0$ and any integer $n\ge0$,
choose $p>0$ so that
\[
 p\le1/100,\quad p\le\epsilon/(Z+\epsilon),\quad
 pC_0\le\epsilon/100,\quad z_p\le Z_0\,2^{-n-3}.
 \tag{B}\label{eq:scalar-choice}
\]
Such a choice exists. If
$(1-\epsilon/4)z_p\le z\le z_p$ and $L=\lfloor Z/z\rfloor$, with
$Z\ge3Z_0/4$, then $L\ge1$ and, for
$Z'=Lz$, $\epsilon'=1-(1-p\epsilon)^L$,
\[
 0\le Z-Z'<z,\qquad
 \frac1{10}<Lp\le1-\frac{3\epsilon}{4},\qquad
 \frac{\epsilon}{12}\le\epsilon'\le\epsilon-\frac{3\epsilon^2}{4}.
 \tag{C}\label{eq:scalar-step-numbers}
\]
\end{lemma}
\begin{proof}
Formula (A) is continuous in $p$ and tends to zero as $p\downarrow0$;
every other constraint in (B) holds for sufficiently small positive $p$.
The condition $p\le\epsilon/(Z+\epsilon)$ gives
\[
 \frac{pZ}{1-\epsilon}\le z_p
 \le\frac{p(Z+\epsilon^2)}{1-\epsilon}.
\]
Since $z\le Z_0/8$ and $Z\ge3Z_0/4$, $L\ge1$.
The floor inequalities give $0\le Z-Lz<z$ and
\[
 \frac{(1-\epsilon)Z}{Z+\epsilon^2}-p<Lp
 \le\frac{1-\epsilon}{1-\epsilon/4}\le1-3\epsilon/4.
\]
The leftmost expression is at least
$(1-1/400)\,3/19-1/100=1121/7600>1/10$.
For $0<t<1$ and $L\ge1$,
$Lt/(1+Lt)\le1-(1-t)^L\le Lt$:
the upper bound is Bernoulli's inequality, and the lower follows
from $(1-t)^{-L}\ge1+Lt$. Apply this with $t=p\epsilon$.
The lower estimate exceeds $\epsilon/(10(1+\epsilon))\ge\epsilon/12$;
the upper estimate is at most $\epsilon(1-3\epsilon/4)$.
\end{proof}

\begin{lemma}[Preserving inverse estimates]\label{lem:scalar-inverse-step}
Take scalar data with loss $\epsilon$ and gain $Z$ satisfying
Lemma~\ref{lem:scalar-numerics}, and suppose $Z_0\le a^2/3$.
Suppose their augmented Gram matrix satisfies
$B\ge\epsilon^2I$ and
$(B^{-1})_{kk}\le1+2\epsilon$ for $k\ne0$.
Take $C_0=\max_j((B^{-1})_{jj}-1)$, including $j=0$.
Activate, choose the finite product family with gain $z$, and take
$L$ independent copies as above. Their augmented matrix $B'$ satisfies
\[
 B'\ge(\epsilon')^2I,\qquad
 ((B')^{-1})_{kk}\le1+2\epsilon'\quad(k\ne0).
\]
\end{lemma}
\begin{proof}
Put $M=1-p\epsilon$ and $h=\|B^{-1}\|\le\epsilon^{-2}$.
If $A_W$ is the one-copy unnormalized centered matrix, then the
new normalized centered inverse is
\[
 D_0=M\diag(A_W^{-1},\ldots,A_W^{-1})\le MhI.
\]
Its diagonal entries are at most $M(1+pC_0)$ by
Lemma~\ref{lem:rare-data}. With new mean column $\mathbf a$,
$\mathbf a^*D_0\mathbf a=Z'\le Z_0$, and
Lemma~\ref{lem:scalar-identities} gives
\[
 (1-\epsilon')(B')^{-1}
 =\begin{pmatrix}1+Z'&-\mathbf a^*D_0\\-D_0\mathbf a&D_0\end{pmatrix}.
\]
For a nonconstant diagonal entry, subtracting one gives the bound
\[
 \frac{\epsilon'+\epsilon/100}{1-\epsilon'}
 \le\frac{\epsilon'+12\epsilon'/100}{9/10}
 =\frac{56}{45}\epsilon'<2\epsilon'.
\]
Here $\epsilon'\le a<1/10$ and $\epsilon\le12\epsilon'$.
For the operator norm use Cauchy--Schwarz in the $D_0$ metric.
The quadratic form of the displayed matrix is bounded by the
two-by-two form with diagonal $1+Z',Mh$ and off-diagonal
$\sqrt{MhZ'}$. Its largest eigenvalue is at most its trace, whence
$\|(B')^{-1}\|\le(h+1+Z_0)/(1-\epsilon')$.
The function $(1-x)/x^2$ decreases on $(0,1)$.
Using (C) and subtracting explicitly,
\[
 \frac{1-\epsilon'}{(\epsilon')^2}-\frac1{\epsilon^2}
 \ge\frac{1/2+3\epsilon/16}{\epsilon(1-3\epsilon/4)^2}
 \ge\frac1{2\epsilon}>1+Z_0.
\]
The last inequality uses $\epsilon\le1/400$ and $Z_0\le a^2/3$.
This bounds the inverse norm by $(\epsilon')^{-2}$, proving the floor.
\end{proof}

\begin{lemma}[Preserving prefix estimates]\label{lem:scalar-prefix-step}
Let $0<\delta\le1$ and $a=\delta/400$.
Suppose the old scalar family is an independent-copy concatenation
with centered bound $\sqrt Q$, total gain $Z$, and each block's gain
at most $\eta$. Suppose
\[
 1+\delta/4\le Q<1+\delta/2,\quad Z\le Z_0\le a^2/3,\quad
 0<\epsilon\le a,
\]
and its nonconstant inverse diagonals are at most $1+2\epsilon$.
The activation and copying step above produces a centered bound
$\sqrt{Q'}$, where
$Q'=Q+20\eta/\delta+\epsilon^2/2$.
The old augmented squared prefix norms are at most $Q+20\eta/\delta$.
\end{lemma}
\begin{proof}
Each centered block inherits the bound $\sqrt Q$ by restriction.
$Q-1-Z>\delta/8$ and $Q^2\le(1+\delta/2)^2\le9/4$, so
$Q^2/(Q-1-Z)<20/\delta$. Lemma~\ref{lem:join} bounds the old
augmented prefixes. For a nonconstant activated row,
Lemma~\ref{lem:row-split} adds at most $\epsilon^2/2$.
For row zero its sharper bound, using the first-prefix identity, is
\[
 1+Z+\frac{\epsilon(Z+\epsilon)}{4(1-\epsilon)}
 \le1+a^2<Q.
\]
Indeed the increment above one is at most
$a^2[1/3+(1+a/3)/(4(1-a))]\le2a^2/3<a^2$, since $a\le1/5$.
Centering preserves each tail-prefix bound, and independent centered
copies are orthogonal. Apply
Lemmas~\ref{lem:prefix-operations} and \ref{lem:scalar-copies}.
\end{proof}

\begin{proposition}[Measurable scalar supply]\label{prop:raw-scalar}
Fix $0<\delta\le1$ and $a,Z_0$ as in
Lemma~\ref{lem:scalar-initial}. There exist scalar data indexed by
$n\in\N$ such that
\[
 \epsilon_n\downarrow0,\quad 3Z_0/4\le Z_n\le Z_0,\quad
 B_n\ge\epsilon_n^2I,\quad
 (B_n^{-1})_{kk}\le1+2\epsilon_n\ (k\ne0).
\]
Both augmented and centered families have a prefix bound strictly
smaller than $1+\delta$.
\end{proposition}
\begin{proof}
Start with Lemma~\ref{lem:scalar-initial}, $Q_0=1+\delta/4$,
and one block of gain $\eta_0=Z_0$. At step $n$ use (B), choose
the finite truncation by Lemma~\ref{lem:rare-data}, and take the
number of copies specified in Lemma~\ref{lem:scalar-numerics}.
Set $\eta_{n+1}=z$ and
$Q_{n+1}=Q_n+20\eta_n/\delta+\epsilon_n^2/2$.
The copy structure is retained, and the preceding three lemmas give
all inverse and prefix assertions at the next stage.

We verify that their hypotheses remain available throughout.
The gain loss is at most $Z_0\,2^{-n-3}$, so
$0\le Z_0-Z_n\le Z_0(1-2^{-n})/4$.
Also $\epsilon_{n+1}\le\epsilon_n-3\epsilon_n^2/4$ and
$\epsilon_{n+1}\ge\epsilon_n/12>0$.
Summing gives $\sum_n\epsilon_n^2\le4a/3$; monotonicity then forces
$\epsilon_n\to0$. Further,
$\sum_n\eta_n\le5Z_0/4<3Z_0/2$.
Consequently
\[
 Q_n\le1+\delta/4+30Z_0/\delta+2a/3
 \le1+\delta/4+\delta/16000+\delta/600<1+\delta/2.
\]
This closes the induction on $Q_n$.
The augmented squared bound is at most
$Q_n+20\eta_n/\delta<1+\delta/2+\delta/24000<(1+\delta)^2$;
the centered squared bound is at most $Q_n$.
Every step uses finitely many coordinates and retains a positive
pointwise lower bound on its weight. No common lower bound across
all stages is asserted.
\end{proof}

\subsection{Finite smoothing and exact gain}
\begin{lemma}[The finite Fejer kernel]\label{lem:fejer}
On $\T$ define
$K_N(y)=(N+1)^{-1}|\sum_{j=0}^N e^{2\pi ijy}|^2$.
It is nonnegative, even, has integral one, and equals
\[
 K_N(y)=\sum_{|l|\le N}\left(1-\frac{|l|}{N+1}\right)e^{2\pi ily}.
\]
On $\T^I$ use the product of these kernels. Its Fourier multiplier
at $k\in\Z^I$ is
$\gamma_N(k)=\prod_{i\in I}(1-|k_i|/(N+1))_+$,
which tends to one for each fixed $k$.
\end{lemma}
\begin{proof}
Expand the squared sum. For each difference $l=j-j'$ there are
$N+1-|l|$ pairs if $|l|\le N$ and none otherwise.
Haar orthogonality gives the integral, and the coefficients prove
evenness. Positivity follows from the defining square. Products,
Fubini, and the finite number of coordinates prove all remaining claims.
\end{proof}

\begin{lemma}[Smoothing one fixed scalar datum]\label{lem:smoothing}
For scalar data on $\T^I$ set $w_N=w*K_N$. Its mass equals that of $w$,
it is an even character polynomial with the same positive lower bound
and upper bound one, and its augmented Gram matrix satisfies
\[
 B_N=\int U(y)^*BU(y)K_N(y)\,dy,\qquad
 U(y)=\diag(1,\chi_1(y),\ldots,\chi_d(y)).
\]
Every augmented prefix bound and every lower bound $B\ge aI$
are preserved. Also $(B_N^{-1})_{kk}\le(B^{-1})_{kk}$ and
$B_N\to B$, so the gains converge.
\end{lemma}
\begin{proof}
Positivity, the lower bound and the upper bound follow because convolution
is a probability average. Translation invariance and Fubini preserve
mass. Expanding the finite kernel and translating the integration
variable writes $w_N$ as a finite character sum. Inversion of all
coordinates proves evenness.
The $(i,j)$ entry of $B_N$ is the corresponding entry of $B$
multiplied by $\gamma_N(k_j-k_i)$, which proves the averaging formula
and entrywise convergence.
Every coordinate projection commutes with $U(y)$, so averaging
$D_kBD_k\le K^2B$ preserves that inequality, and also preserves $aI$.
Lemma~\ref{lem:inverse-variational} gives
$B_N^{-1}\le\int U(y)^*B^{-1}U(y)K_N(y)\,dy$;
its diagonal is that of $B^{-1}$.
Inverse continuity and the fixed mass in
Lemma~\ref{lem:scalar-identities} prove convergence of gain.
The centered prefix bound follows by projecting off the initial constant.
\end{proof}

\begin{lemma}[Tuning a scalar gain]\label{lem:tuning}
Suppose polynomial scalar data have loss $\epsilon>0$, gain $Z>z>0$,
augmented bound $K\ge1$, and
$B\ge\epsilon^2I$, $(B^{-1})_{kk}\le1+2\epsilon$ for $k\ne0$.
There are polynomial scalar data on the same characters with gain
exactly $z$, loss $0<\epsilon'\le\epsilon$, the same bound $K$ for
augmented and centered families, and the corresponding inverse estimates.
\end{lemma}
\begin{proof}
For $0\le t\le1$ put $w_t=1-t+tw$, so
$B_t=(1-t)I+tB$ and $\epsilon_t=t\epsilon$.
This remains a positive polynomial weight. Since $K\ge1$,
the prefix inequality holds for the convex combination of $I$ and $B$.
Also $B_t\ge(1-t+t\epsilon^2)I\ge(t\epsilon)^2I$.
Inverse averaging with two points gives
$(B_t^{-1})_{kk}\le1-t+t(B^{-1})_{kk}\le1+2t\epsilon$.
Its gain is continuous, zero at $t=0$, and $Z>z$ at $t=1$.
The intermediate value theorem gives a solution $t>0$.
Centering is covered by Lemma~\ref{lem:prefix-operations}.
\end{proof}

\begin{proposition}[Polynomial scalar supply]\label{prop:scalar-supply}
For $0<\delta\le1$ let $a=\delta/400$,
$Z_0=a^2/(4-8a+3a^2)$, and $c_\delta=Z_0/2$.
For every $0<z<c_\delta$ and $\eta>0$, there are polynomial scalar data
with gain $z$ and loss $0<\epsilon<\eta$ such that
\[
 \begin{gathered}
 \PB_K(1,\chi),\quad \PB_K(\chi^\circ)
 \quad\hbox{for some }K<1+\delta,\\
 B\ge\epsilon^2I,\quad (B^{-1})_{kk}\le1+2\epsilon\ (k\ne0).
 \end{gathered}
\]
\end{proposition}
\begin{proof}
Choose a datum in Proposition~\ref{prop:raw-scalar} with loss less
than $\eta$. Smooth this fixed finite datum. Its gains tend to a
value at least $3Z_0/4$, so for sufficiently large smoothing index
the gain exceeds $Z_0/2$. Lemma~\ref{lem:smoothing} preserves the
prefix and inverse estimates. Apply Lemma~\ref{lem:tuning} to obtain
the exact gain. Its positive loss does not increase.
\end{proof}

\section{A directional finite construction}
\subsection{Response and scale}
\begin{lemma}[Uniform scalar response estimates]\label{lem:response}
Fix $\delta$ and $0<z<c_\delta$, and consider any datum from
Proposition~\ref{prop:scalar-supply} with loss $\epsilon$.
For $0\le\beta\le1$ put
\[
 \Phi(\beta)=(1-\beta)^2b^*[\beta I+(1-\beta)G]^{-1}b.
\]
This function is nonnegative and decreasing. For all sufficiently small
$\epsilon$, uniformly over these data,
\[
 \Phi(\epsilon^2)\ge\frac{z}{8(1+z)},\qquad
 \Phi(\sqrt\epsilon)\le\sqrt\epsilon,\qquad
 \Phi(R\beta)\ge\frac{\Phi(\beta)}{4R}
 \quad(R\ge1,\ 0\le R\beta\le1/2).
\]
\end{lemma}
\begin{proof}
Diagonalize $0<G\le I$. Each summand of $\Phi$ is a nonnegative
multiple of $(1-\beta)^2/(\beta+(1-\beta)g)$, with $0<g\le1$;
its derivative is nonpositive. The scalar projection identity gives
$\Phi(0)=(1-\epsilon)z/(1+z)$.
Since $G\ge\epsilon^2I$,
$\epsilon^2I+(1-\epsilon^2)G\le2G$, so, for $\epsilon\le1/2$,
the first bound follows by inverse order and
$(1-\epsilon^2)^2\ge1/2$.
For each nonconstant character,
$b_k=-\int\overline{\chi_k}(1-w)$. The finite orthonormal projection
inequality gives $\|b\|^2\le\int(1-w)^2\le\epsilon$.
The inverse at $\beta=\sqrt\epsilon$ is at most
$\epsilon^{-1/2}I$, proving the second bound.
Finally
$R\beta I+(1-R\beta)G\le R[\beta I+(1-\beta)G]$ and
$(1-R\beta)^2\ge(1-\beta)^2/4$ prove the ratio estimate.
\end{proof}

\begin{lemma}[A geometric-grid scale]\label{lem:scale}
Let $0<\epsilon<1$, $R\ge2$, and
$\log(1/\epsilon)\ge8\log(4R)$.
Suppose $\Phi:[0,1]\to[0,\infty)$ decreases and satisfies
$\Phi(\epsilon^2)>\epsilon^{1/4}$,
$\Phi(\sqrt\epsilon)\le\sqrt\epsilon$, and
$\Phi(R\beta)\ge\Phi(\beta)/(4R)$ when $R\beta\le1/2$.
Then some $m\in[\epsilon^2R,R\sqrt\epsilon]$ satisfies
\[
 \sqrt\epsilon<\Phi(m/R)\le\epsilon^{1/4},\qquad
 \Phi(m)\le R^{-1/16}\Phi(m/R).
\]
\end{lemma}
\begin{proof}
Take $x_k=\epsilon^2R^k$ until the first $K$ with
$x_K\ge\sqrt\epsilon$. Then
$x_K\le R\sqrt\epsilon\le1/2$ and
$K\le3\log(1/\epsilon)/(2\log R)+1$.
All values $\Phi(x_k)$ are positive, by the ratio bound and induction.
Let $a,b$ be the first crossings of $\epsilon^{1/4},\epsilon^{1/2}$.
They exist; $a\ge1$ and
$\Phi(x_a)>\epsilon^{1/4}/(4R)\ge\epsilon^{1/2}$, so $a<b$.
If every step $a<k\le b$ had
$\Phi(x_k)>R^{-1/16}\Phi(x_{k-1})$, then
\[
 \tfrac14\log(1/\epsilon)-\log(4R)
 <\log\frac{\Phi(x_a)}{\Phi(x_b)}
 <\tfrac3{32}\log(1/\epsilon)+\tfrac1{16}\log R.
\]
This is impossible: subtracting the rightmost bound from the leftmost,
and using the logarithmic hypothesis, gives at least
$\frac14\log(4R)-\frac1{16}\log R>0$.
Choose a step where the inequality fails. Its $m=x_k$ has the desired
bounds, since $a\le k-1<b$.
\end{proof}

\subsection{Seed data and a uniform parameter sequence}
For the rest of this section fix a probability $\rho$ on $X_I$, an
ordered independent character family $(\phi_j)_{j=1}^d$ with bound
$K_V\ge1$, its span $V$, and a unit vector $e=\sum_jc_j\phi_j$.
Choose
\[
 \kappa=\bigl(4\max(1,\sum_j|c_j|)\bigr)^{-1},\qquad
 \theta=\kappa^2,\qquad a=1-\theta.
\]
These quantities remain fixed while the scalar data change.
In particular $0<\theta\le1/16$ and $|\kappa e(s)|\le1/4$ at every
base point, not just almost everywhere.

\begin{definition}[The seed]\label{def:seed}
For an integer $L\ge1$ and $0<m\le2^{-L}$, introduce phase coordinates
$y_1,\ldots,y_L$ and put
\[
 \zeta_\ell(y)=e^{2\pi iy_\ell},\qquad
 w_e(s,y)=m\prod_{\ell=1}^L
 [1+2\kappa\operatorname{Re}(e(s)\overline{\zeta_\ell(y)})].
\]
The reference measure for rows is $\rho\,dy$.
Let $B_e$ be the compression of multiplication by $w_e$ to
$\mathcal R=V\oplus\spn(\zeta_1,\ldots,\zeta_L)$ in that reference
Hilbert space.
\end{definition}

\begin{lemma}[Seed moments and spectral data]\label{lem:seed-spectrum}
The seed is a character polynomial, $0<w_e\le1$, and its phase
integral is the constant $m$. In particular $0<B_e\le I$.
Under the normalized seed fibers,
\[
 \E[\zeta_\ell\mid s]=\kappa e(s),\qquad
 \E[\overline{\zeta_\ell}\zeta_k\mid s]=\theta|e(s)|^2
 \quad(\ell\ne k).
\]
The operator $B_e$ is $mI$ on $V\cap e^\perp$, $maI$ on the
span of the differences of the $\zeta_\ell$, and has matrix
\[
 mH_L=m\begin{pmatrix}1&\sqrt{L\theta}\\
 \sqrt{L\theta}&a+L\theta\end{pmatrix}
\]
on the reference orthonormal pair $e,L^{-1/2}\sum_\ell\zeta_\ell$.
Its eigenvalues there are $m\gamma,m\Gamma$, with
\[
 0<\gamma<1<\Gamma,\quad \gamma\Gamma=a,\quad
 1+\frac{\theta L}{a}\le R:=\gamma^{-1}
 \le\frac{2+\theta L}{a}.
\]
The squared norm of the projection of $e$ onto the $\gamma$ eigenspace
of $H_L$ is
$p_\gamma=(\Gamma-1)/(\Gamma-\gamma)$ and tends to one as $L\to\infty$.
\end{lemma}
\begin{proof}
Each factor lies in $[1/2,3/2]$ and has phase integral one.
Thus $m2^{-L}\le w_e\le m(3/2)^L\le1$; finite polynomiality follows
from Lemma~\ref{lem:character-algebra}. Integrating one character
against its own factor gives $\kappa e(s)$; different factors are
independent in the phase variables, giving the mixed moments.
The $\zeta_\ell$ are orthonormal and orthogonal to $V$ in the reference
measure. The moments just computed, together with $\|e\|_\rho=1$,
give the stated orthogonal decomposition and matrix.
The determinant and trace of $H_L$ are $a$ and $1+a+L\theta$.
Also $\det(H_L-I)=-L\theta<0$, so its two positive eigenvalues
straddle one. By its diagonal Rayleigh quotient and its trace,
$a+L\theta\le\Gamma\le2+L\theta$; use $R=\Gamma/a$.
The spectral projection is
$(\Gamma I-H_L)/(\Gamma-\gamma)$, proving the formula for $p_\gamma$.
For $L>1$,
$0\le1-p_\gamma\le1/(\Gamma-1)\le1/(\theta(L-1))$.
This proves the asserted limit.
\end{proof}

\begin{lemma}[Compatible scalar losses and seed scales]\label{lem:seed-scales}
Fix $K_s>1$, choose $0<\delta\le1$ with $1+\delta<K_s$, and fix
$0<z<c_\delta$. There are scalar data with losses $\epsilon_n\to0$,
response functions $\Phi_n$, and seed parameters $L_n,m_n$ such that,
after omitting finitely many indices,
\[
 L_n=\left\lfloor\frac{\log(1/\epsilon_n)}{4\log2}\right\rfloor\ge1,
 \quad 0<m_n\le2^{-L_n},\quad p_{\gamma_n}\ge3/4,
\]
\[
 \sqrt{\epsilon_n}<\Phi_{*,n}:=\Phi_n(m_n\gamma_n)
 \le\epsilon_n^{1/4},\quad
 \frac{\Phi_n(m_n)}{\Phi_{*,n}}\le R_n^{-1/16}\to0.
\]
Each scalar augmented family has bound less than $K_s$.
\end{lemma}
\begin{proof}
Choose the $n$th scalar datum with exact gain $z$ and loss less than
$2^{-n-1}$ using Proposition~\ref{prop:scalar-supply}.
Put $t_n=\log(1/\epsilon_n)\to\infty$. Lemma~\ref{lem:seed-spectrum}
gives $R_n\to\infty$ and $R_n\le C(1+t_n)$ for a fixed $C$.
The elementary estimates
$\log y\le2\sqrt y$ for $y\ge1$ and
$e^{t/8}\ge t^2/128$ show
\[
 \frac{\log(4C(1+t))}{t}\to0,\qquad
 C(1+t)e^{-t/4}\to0.
\]
The logarithm estimate follows by integrating
$1/s\le1/\sqrt s$ on $[1,y]$; the exponential estimate is its
nonnegative quadratic power-series term for $t\ge0$.
For the first limit, divide the logarithm estimate by $t$; for the second,
bound $t e^{-t/4}$ by $128t^{-1}e^{-t/8}$ for $t>0$.
Thus eventually $R_n\ge2$,
$t_n\ge8\log(4R_n)$, and
$R_n\sqrt{\epsilon_n}<\epsilon_n^{1/4}\le2^{-L_n}$.
Also $p_{\gamma_n}\ge3/4$ eventually.
The fixed positive lower bound in Lemma~\ref{lem:response} eventually
exceeds $\epsilon_n^{1/4}$. Apply Lemma~\ref{lem:scale} to each
individual $\Phi_n,R_n,\epsilon_n$ to choose $m_n$.
This uses uniform inequalities; it asserts no continuity of a scalar
packet, its dimension, or its matrix as a function of the loss.
\end{proof}

\subsection{One selected family}
\begin{definition}[The selected family]\label{def:selected}
Use one of the scalar systems above, with scalar variables disjoint
from the phase and base variables. Put
\[
 W=w+(1-w)w_e,\qquad M=1-(1-m)\epsilon.
\]
For each nonconstant scalar character $\chi_k$, list the rows
\[
 \phi_1\chi_k,\ldots,\phi_d\chi_k,
 \zeta_1\chi_k,\ldots,\zeta_L\chi_k,
\]
in that order, and order the columns by $k$. This is the new family.
\end{definition}

\begin{lemma}[Selected means and independence]\label{lem:selected-means}
$W$ is a positive character polynomial at most one, with auxiliary
integral $M$ independent of the base point. Its conditional means are
\[
 \mathcal E(\phi_j\chi_k)=\frac{(1-m)b_k}{M}\phi_j,\qquad
 \mathcal E(\zeta_\ell\chi_k)=-\frac{m\kappa b_k}{M}e.
\]
Its centered Gram matrix is positive definite.
\end{lemma}
\begin{proof}
$W=1-(1-w)(1-w_e)$ lies between $w>0$ and one; its integral follows
by first integrating the phase variables, then the scalar variables.
The scalar coefficients $b_k$ are real by evenness.
Use $\int\chi_k=0$, $\int\chi_kw=b_k$, and the seed means to get both
formulas.
If a centered linear combination vanishes, positivity of $W$ makes
the same equality valid almost everywhere for the reference product
measure. It says that a sum of nonconstant scalar characters with
row coefficients equals a base function. Haar integration in the
scalar variables makes that base function zero. Each scalar Fourier
coefficient therefore vanishes. Averaging in the phase variables
gives a zero linear combination of the independent $\phi_j$, so its
coefficients vanish. The remaining phase Fourier coefficients then
make every $\zeta_\ell$ coefficient zero.
\end{proof}

\begin{lemma}[The two eigenvalues of the information operator]
\label{lem:selected-information}
For the selected family, its information operator on the base is
\[
 J=\lambda_\perp P_{V\cap e^\perp}+\lambda P_{\C e},\qquad
 \lambda_\perp=\frac{\Phi(m)}{M-\Phi(m)},\quad
 \lambda=\frac{\Phi_e}{M-\Phi_e},
\]
where
$\Phi_e=p_\gamma\Phi(m\gamma)+(1-p_\gamma)\Phi(m\Gamma)$.
The denominators are positive. Under the choices of
Lemma~\ref{lem:seed-scales}, $3\Phi_*/4\le\Phi_e\le\Phi_*$.
\end{lemma}
\begin{proof}
In the reference row space diagonalize $B_e$, with orthonormal
eigenvectors $u_i$ and eigenvalues $0<\beta_i\le1$.
From $W=w_e+(1-w_e)w$, the unnormalized Gram matrix of
$(u_i\chi_k)_k$ is $\beta_iI+(1-\beta_i)G$.
Different $i$ are orthogonal; for a base vector $v\in V$, its pairing
column is $(1-\beta_i)b\,\ip{u_i}{v}$.
Lemma~\ref{lem:gram-projection} therefore gives the unnormalized
projection form
$\sum_i\Phi(\beta_i)|\ip{u_i}{v}|^2$.
On $V$ this is $\Phi(m)$ on $V\cap e^\perp$ and $\Phi_e$ on $\C e$,
by Lemma~\ref{lem:seed-spectrum}.
Normalization divides the form by $M$. The means lie in $V$, so
$J$ vanishes on $V^\perp$. Lemma~\ref{lem:information} identifies
the form with $\Psi(J)$, whose eigenvalues are strictly below one
because $J$ is finite and positive. Solving $q=\lambda/(1+\lambda)$
gives the two formulas and positive denominators.
Monotonicity of $\Phi$, $m\Gamma\ge m\gamma$, and
$p_\gamma\ge3/4$ give the last bounds.
\end{proof}

\begin{lemma}[Completed and partial columns]\label{lem:column-bounds}
Adjoin the constant scalar column with all rows to the selected
family, and let $H$ be its unnormalized augmented Gram matrix.
Its reference row Gram matrix is $\Omega=\diag(G_\phi,I_L)$.
Prefixes consisting of whole columns have norm at most $K_s$.
For a nonconstant column injection $I_k$,
\[
 I_k^*HI_k\le\Omega,\qquad
 I_k^*H^{-1}I_k\le(1+2\epsilon)\Omega^{-1}.
\]
Every prefix ending inside a nonconstant column has squared norm
at most
\[
 Q(\epsilon)=1+
 \left(\sqrt{K_s^2-1}+
 \sqrt{(1+2\epsilon)K_V^2-1}\right)^2.
\]
\end{lemma}
\begin{proof}
The augmented family is independent: it is independent in the
reference product measure by scalar Haar orthogonality and row
independence, and $W$ is bounded below by a positive constant.
In orthonormal row eigencoordinates of $B_e$ its scalar blocks are
$\beta I+(1-\beta)B$. Their coordinate-prefix inequalities follow
by taking convex combinations of those for $I$ and $B$.
This bounds completed columns.
The same coordinates give an upper bound one on a column synthesis
Gram block, because $W\le1$. Inverse convexity from
Lemma~\ref{lem:inverse-variational} gives
\[
 [\beta I+(1-\beta)B]^{-1}
 \le\beta I+(1-\beta)B^{-1}.
\]
Its nonconstant diagonal entries are at most $1+2\epsilon$.
To return to the original row coefficients, let $R$ be the
invertible row-synthesis matrix in reference orthonormal coordinates,
so $R^*R=\Omega$. The full Gram transforms as
$H=E^*\widetilde HE$, where $E$ is diagonal with copies of $R$.
Its inverse transforms as
$H^{-1}=E^{-1}\widetilde H^{-1}E^{-*}$.
These formulas give precisely the two displayed column bounds.

The first inequality bounds column synthesis, from the $\Omega$
coefficient norm, by one. The second bounds coefficient extraction
into that norm by $\sqrt{1+2\epsilon}$: for any total coefficient
vector $x$ and row test vector $z$, Cauchy--Schwarz in the $H$ metric
gives
$|z^*x_k|^2\le(z^*I_k^*H^{-1}I_kz)(x^*Hx)$;
take the dual supremum in the $\Omega$ metric.
The initial row cut has norm at most $K_V$, since the phase rows are
orthonormal and orthogonal to $V$ in the reference metric.
Thus a prefix is $P+Q$, where $PQ=QP=0$,
$\|P\|\le K_s$, and $\|Q\|\le\sqrt{1+2\epsilon}K_V$.
It is an idempotent. Lemma~\ref{lem:oblique} and the triangle
inequality for $(P-P^*)+(Q-Q^*)$ prove the bound.
If $P$ or $Q$ is zero, the bound follows from the other one's norm;
if both are zero it is immediate.
\end{proof}

\begin{lemma}[Centering the selected family]\label{lem:selected-prefix}
The centered new family has prefix bound $\sqrt{Q(\epsilon)}$.
\end{lemma}
\begin{proof}
Remove the phase rows from the constant scalar column of the augmented
family. Restriction preserves every prefix bound ending in a
nonconstant column. The remaining initial column spans $V$.
By Lemma~\ref{lem:selected-means}, each new vector has its mean in
$V$, and its centered part is orthogonal to all base functions.
Consequently projection off $V$ is exactly centering.
Apply Lemma~\ref{lem:prefix-operations}(3), which requires bounds
only for prefixes ending after that initial column.
\end{proof}

\begin{lemma}[Vanishing leakage and amplification]\label{lem:amplification}
Along the sequence of Lemma~\ref{lem:seed-scales},
\[
 \lambda_n\to0,\qquad
 \lambda_{\perp,n}/\lambda_n\to0,\qquad
 (1-M_n)/\lambda_n\to0,\qquad \lambda_n>0.
\]
For any $\sigma>0$, with $N_n=\lceil\sigma/\lambda_n\rceil$,
the product of $N_n$ independent selected copies has mass tending
to one and information operator tending in norm to $\sigma P_{\C e}$.
Its centered squared bound remains $Q(\epsilon_n)$.
\end{lemma}
\begin{proof}
Eventually $M_n\ge3/4$ and $\Phi_{*,n}\le1/4$. Thus
\[
 \tfrac34\Phi_{*,n}\le\lambda_n\le2\Phi_{*,n},\quad
 \frac{\lambda_{\perp,n}}{\lambda_n}\le
 \tfrac83 R_n^{-1/16},\quad
 \frac{1-M_n}{\lambda_n}\le\tfrac43\sqrt{\epsilon_n}.
\]
The preceding lemmas give all these inequalities, including positivity.
Now $\sigma\le N_n\lambda_n<\sigma+\lambda_n$, whence
$N_n\lambda_{\perp,n}\to0$ and $N_n(1-M_n)\to0$.
Lemma~\ref{lem:independent-blocks} adds the information operators and
preserves centered bounds. The product mass satisfies
$1-M_n^{N_n}\le N_n(1-M_n)$.
\end{proof}

\begin{proposition}[Directional supply]\label{prop:directional}
For the fixed $V,e$, any $K_s>1$, $\sigma>0$ and $\eta>0$, there
are a positive character-polynomial weight $W\le1$ with constant
new-variable integral $M>1-\eta$, and an ordered character family.
Its means, centered Gram matrix, and information operator are computed
in $L^2(M^{-1}W\,d\rho\,dy)$. The means lie in $V$, the centered
Gram matrix is positive definite, and
\[
 \|J-\sigma P_{\C e}\|<\eta,\qquad
 \PB_{\sqrt{1+(\sqrt{K_s^2-1}+\sqrt{K_V^2-1})^2+\eta}}
       (\hbox{centered family}).
\]
Every character has a nonzero frequency in the new variables.
Its real frequency is one of those of the $\phi_j$, or zero.
If the $\phi_j$ have pairwise distinct old auxiliary modes, the
output has pairwise distinct full auxiliary modes.
\end{proposition}
\begin{proof}
Use the sequence of Lemma~\ref{lem:seed-scales} and the amplification
of Lemma~\ref{lem:amplification}.
Since
\[
 Q(\epsilon_n)\longrightarrow
 1+(\sqrt{K_s^2-1}+\sqrt{K_V^2-1})^2,
\]
one sufficiently large index gives all three required tolerances
simultaneously.
The finite product weight remains positive, at most one, and polynomial.
The means remain in $V$ by integration of the unused factors.
Within one copy every column has a nonconstant scalar mode.
Helper rows retain their old modes; phase rows have their separate
phase-coordinate unit modes. Therefore distinct old auxiliary modes
distinguish helper rows, and phase modes distinguish the other rows.
Different copies have disjoint scalar blocks containing nonzero modes,
so they cannot collide. The real-frequency assertion follows from the
explicit products in Definition~\ref{def:selected}.
\end{proof}

\section{Real exponentials and the basis interface}
\subsection{Totality and independence}
\begin{lemma}[Continuity of a density's Fourier transform]
\label{lem:density-transform}
If $\nu$ is a finite positive Borel measure on $\R$ and $h\in L^2(\nu)$,
then $h\in L^1(\nu)$ and
$F_h(t)=\int h(x)e_{-t}(x)\,d\nu(x)$ is continuous.
If $F_h=0$, then $h=0$ in $L^2(\nu)$.
\end{lemma}
\begin{proof}
Cauchy--Schwarz gives $\|h\|_1\le\sqrt{\nu(\R)}\|h\|_2$.
Dominated convergence applies with majorant $|h|$.
For $a=\operatorname{Re}h$ and $b=\operatorname{Im}h$,
\[
 F_a(t)=\tfrac12(F_h(t)+\overline{F_h(-t)}),\qquad
 F_b(t)=\tfrac1{2i}(F_h(t)-\overline{F_h(-t)}).
\]
Thus both transforms vanish. The finite positive measures $a_+\nu$
and $a_-\nu$ have equal characteristic functions, so Mathlib's
finite-measure uniqueness theorem makes them equal. On $\{a>0\}$
their equality gives $\int a_+\,d\nu=0$; the negative part is treated
on $\{a<0\}$. Hence $a=0$ almost everywhere. The same argument for
$b$ proves the claim. The change from $e_{-t}$ to the characteristic
function convention is the bijective parameter change $t\mapsto-2\pi t$.
\end{proof}

\begin{lemma}[Rational totality]\label{lem:real-totality}
For every finite positive Borel $\nu$ on $\R$,
$\overline{\spn}\{e_q:q\in\Q\}=L^2(\nu)$.
\end{lemma}
\begin{proof}
A vector orthogonal to all these characters has $F_h(q)=0$ for all
rational $q$. Lemma~\ref{lem:density-transform} and density of $\Q$
give $F_h=0$ and then $h=0$. Use the orthogonal-complement density
criterion.
\end{proof}

\begin{lemma}[Pointwise character independence]\label{lem:pointwise-independence}
For distinct real $t_1,\ldots,t_d$, the functions $e_{t_j}$ are
linearly independent as functions on $\R$.
\end{lemma}
\begin{proof}
Each is a monoid character of the additive real group with values in
the multiplicative monoid of $\C$. If $s\ne t$, evaluation at
$x=1/(2(s-t))$ gives $e_s(x)/e_t(x)=-1$, so the characters differ.
Apply Mathlib's independence theorem for distinct monoid homomorphisms,
using the multiplicative type synonym for the additive domain.
This theorem imposes no finiteness assumption on the domain group.
\end{proof}

\begin{lemma}[Independence for a singleton-null measure]
\label{lem:real-independence}
If $\nu\ne0$ is finite and singleton-null, distinct real exponentials
are linearly independent in $L^2(\nu)$.
\end{lemma}
\begin{proof}
Suppose a nontrivial finite combination $p$ were zero almost everywhere.
Its measurable zero set has positive measure. The Mathlib accumulation
point theorem for singleton-null measures supplies an accumulation
point of that set; its separability hypothesis holds because it is
a subset of $\R$. The function $p:\R\to\C$ is real analytic: exponential
is analytic, its argument here is real linear, and the sum is finite.
The analytic identity principle on the connected real line makes
$p$ identically zero. This contradicts
Lemma~\ref{lem:pointwise-independence}.
\end{proof}

\subsection{Continuous coordinates and prefix construction}
\begin{lemma}[The coefficient sequence space]\label{lem:coefficient-space}
Let $(b_n)$ be independent in a Banach space $E$. The space
\[
 \mathcal A=\{a\in\C^\N:\sum_{n<N}a_nb_n
                     \hbox{ converges in }E\},\qquad
 \|a\|_{\mathcal A}=\sup_N\left\|\sum_{n<N}a_nb_n\right\|
\]
is a Banach space. The sum map $\mathcal A\to E$ has norm at most one.
\end{lemma}
\begin{proof}
Convergence bounds the supremum; independence makes it a norm.
For every $n$, $|a_n|\|b_n\|\le2\|a\|_{\mathcal A}$.
If $a^{(j)}$ is Cauchy, each coordinate converges to a scalar $a_n$.
For every $\epsilon>0$ and all sufficiently large $j$, passing to
the coordinate limit in each finite sum gives
$\sup_N\|\sum_{n<N}(a_n-a_n^{(j)})b_n\|\le\epsilon$.
Tail differences of partial sums of $a$ are therefore at most
$2\epsilon$ plus the tail difference of the convergent series for
a fixed such $a^{(j)}$. Completeness of $E$ gives convergence of the
series for $a$; the same supremum bound proves convergence in
$\mathcal A$. The sum-map bound follows by taking limits of partial sums.
\end{proof}

\begin{lemma}[The expansion-only definition has continuous coordinates]
\label{lem:coordinate-continuity}
A Schauder basis in the definition of Section~1 has continuous linear
coordinate functionals, and its partial-sum operators have one common
bound.
\end{lemma}
\begin{proof}
Uniqueness of the expansion of zero gives independence and nonzero
basis vectors. The sum map in Lemma~\ref{lem:coefficient-space} is a
continuous linear bijection onto $E$. Its inverse is continuous by
the bounded inverse theorem. Thus
$\sup_N\|\sum_{n<N}a_n(x)b_n\|\le K\|x\|$.
The coordinate estimate in that lemma gives continuity of every $a_n$;
linearity follows from uniqueness.
\end{proof}

\begin{remark}
Lemma~\ref{lem:coordinate-continuity} preserves the mathematical
definition in the original statement. For the supplied Lean target,
these continuous coordinates are already fields of
\ml{SchauderBasis}; its sequence-space bridge need not be implemented
a second time.
\end{remark}

\begin{lemma}[Weak nullity of a bounded Hilbert basis]\label{lem:weak-null}
If $(b_n)$ is a bounded Schauder basis of a Hilbert space, then
$\ip{y}{b_n}\to0$ for every $y$.
\end{lemma}
\begin{proof}
Represent the continuous coordinates by $a_n(x)=\ip{y_n}{x}$ using
Hilbert duality. A vector orthogonal to all $y_n$ has all coordinates
zero and is therefore zero. Thus their span is dense.
If $y'$ is a finite combination of the first $N$ coordinate vectors,
then $\ip{y'}{b_n}=0$ for $n\ge N$ by biorthogonality.
For $M=\sup_n\|b_n\|$, the remaining pairing is bounded by
$M\|y-y'\|$. Approximate $y$ arbitrarily closely.
\end{proof}

\begin{lemma}[The bounded ordered-prefix criterion]\label{lem:basis-criterion}
Let $(b_n)$ be independent with dense span in a Hilbert space.
Suppose every finite ordered initial family has prefix bound $K$,
where $K\ge1$. Then $(b_n)$ is a Schauder basis.
Its prefix operators $P_N$ satisfy
\[
 P_0=0,\quad \dim\ran P_N=N,\quad
 P_NP_M=P_{\min(N,M)},\quad P_Nx\to x.
\]
The difference $P_{n+1}-P_n$ has range $\C b_n$.
\end{lemma}
\begin{proof}
On the algebraic span define $P_N$ by truncating the unique finite
coefficient representation. Its norm is at most $K$, independently of
which larger initial family contains the representation.
For $p_j\to x$ in that span, define $P_Nx=\lim_jP_Np_j$.
The bound makes the images Cauchy and the limit independent of
approximants. Linearity and the bound persist.
The range is the span of the first $N$ vectors: containment follows
from its finite-dimensional closedness, and equality follows because
those vectors are fixed. Independence gives its dimension.
The composition identities hold on the dense span and therefore
everywhere. If $p$ is a finite sum and $N$ contains it,
$\|P_Nx-x\|\le(K+1)\|x-p\|$, proving convergence.
The same density argument gives the stated difference range; it
contains $b_n$ because $(P_{n+1}-P_n)b_n=b_n$.

Define $a_n(x)$ by $(P_{n+1}-P_n)x=a_n(x)b_n$.
The nonzero vector $b_n$ makes this scalar unique and linear, and
$|a_n(x)|\le2K\|x\|/\|b_n\|$ makes it continuous.
Biorthogonality follows by applying the differences to $b_j$.
Telescoping gives $\sum_{n<N}a_n(x)b_n=P_Nx\to x$.
Applying each continuous $a_n$ to any other convergent expansion proves
uniqueness. Equivalently, the displayed projection data and the
specified vectors fulfill Mathlib's
\ml{SchauderBasis.RankOneDecomposition} constructor.
\end{proof}

\begin{lemma}[Transport of a basis]\label{lem:basis-transport}
A linear isometry equivalence sends a Schauder basis to a Schauder
basis, preserving every prefix bound.
\end{lemma}
\begin{proof}
Apply the equivalence to the convergent expansion of its inverse image.
The inverse makes the resulting expansion unique. Isometry preserves
the finite prefix inequalities; the transported coordinates are the
old continuous coordinates composed with the inverse.
\end{proof}

\section{Residual copies and a finite product extension}
Fix $0<u\le1/100$, and set
$K_*=1+10u$ and $c=u/(1+2u)$.
Let $\rho$ be a singleton-null probability on $\R$, let $A$ be a
nonempty ordered list of distinct rational exponentials, and let $f$
be a rational exponential not in that list. Put
$E=\spn A$, $r=f-P_Ef$, and $e=r/\|r\|_\rho$.
Lemma~\ref{lem:real-independence} gives $r\ne0$.

\begin{lemma}[Lifting the residual with prescribed scalar gain]
\label{lem:residual-lift}
Write $\phi=(A,f)$, $r=\sum_ja_j\phi_j$, and let $S$ be the Gram
matrix of $\phi$ under $\rho$. Fix $0<z<c_u$ with $z\|S\|<u$.
For every $a_*>0$ there are a finite product torus, a polynomial
weight $0<Q_*\le1$ depending only on its coordinates with
$\int Q_*>1-a_*$, and an ordered family
$(\phi_j\psi_{jk})_{j,k}$.
In $L^2(d\rho\,Q_*\,dy/\int Q_*)$ this family has bound
$\sqrt{(1+u)^3}$ and has a polynomial $h$ in its span such that
\[
 \E[h\mid x]=r,\qquad
 \|h\|^2=\|r\|_\rho^2+z^{-1}\sum_j|a_j|^2.
\]
The family has pairwise distinct nonzero auxiliary modes.
\end{lemma}
\begin{proof}
For each $j$ take independent scalar data from
Proposition~\ref{prop:scalar-supply}, with parameters $u,z$, on
disjoint torus coordinates. Choose their finitely many losses to
sum to less than $a_*$. Their product $Q_*$ has loss at most this sum.
Use the product of their normalized weights with $\rho$.
The centered blocks of $\phi_j\psi_{jk}$ are orthogonal for distinct
$j$; within a block their Gram is $C_j$, since $|\phi_j|=1$.
The mean synthesis map therefore satisfies
\[
 TC^{-1}T^*=z\sum_j(\phi_j\otimes\phi_j),\qquad
 C=\diag(C_j),
\]
where $(v\otimes v)x=\ip{v}{x}v$.
The norm of $\sum_j\phi_j\otimes\phi_j$ equals $\|S\|$ by
Lemma~\ref{lem:operator-norm} applied to synthesis onto its finite
range. Lemma~\ref{lem:information} bounds the uncentered Gram between
$C$ and $(1+u)C$. The centered families have bounds less than
$1+u$, so Lemma~\ref{lem:prefix-operations} gives the asserted bound.

Choose in each scalar family the mean-one polynomial $g_j$ with
variance $1/z$ from Lemma~\ref{lem:scalar-identities}.
Set $h=\sum_j a_j\phi_jg_j$.
Its mean is $r$. Its centered summands from different $j$ are
orthogonal and have squared norms $|a_j|^2/z$, proving the formula.
The scalar modes are nonzero and distinct within each block; disjoint
blocks distinguish different $j$. The number of blocks is fixed,
so their losses can sum to any prescribed positive bound.
\end{proof}

\begin{lemma}[Equicorrelated unit vectors]\label{lem:equicorrelated}
Suppose $h_i=\alpha e+\xi_i$ for $1\le i\le N$, where
$\|e\|=\|h_i\|=1$, $0<\alpha<1$, and the $\xi_i$ are mutually
orthogonal and orthogonal to $e$. Let $N=\lceil\alpha^{-2}\rceil$.
If each $h_i\perp E$, then
$J_0=u\sum_iP_{\C h_i}$ satisfies
\[
 J_0E=0,\quad \|J_0\|=u[1+(N-1)\alpha^2]<2u,\quad
 \ip{e}{\Psi(J_0)e}\ge c.
\]
\end{lemma}
\begin{proof}
The Gram matrix of the $h_i$ is
$(1-\alpha^2)I+\alpha^2\mathbf1\mathbf1^*$.
Its largest eigenvalue is $1+(N-1)\alpha^2$; synthesis and
Lemma~\ref{lem:operator-norm} give the norm formula.
The ceiling inequalities give $(N-1)\alpha^2<1$ and $N\alpha^2\ge1$.
Also $\ip{e}{J_0e}=uN\alpha^2\ge u$.
On the finite range, diagonalization gives
$\Psi(J_0)\ge J_0/(1+\|J_0\|)$, proving the lower bound by
$c=u/(1+2u)$. On its orthogonal complement both operators vanish.
The annihilation of $E$ follows from the stated orthogonality.
\end{proof}

\begin{lemma}[Residual copies with controlled total loss]
\label{lem:residual-copies}
For every $a>0$ there are an integer $N\ge1$ and a polynomial torus
weight $0<Q\le1$ with
$\int Q>1-a$, and finite ordered character families $V_i$ with unit
vectors $h_i\in\spn V_i$, in
$L^2(d\rho\,Q\,dy/\int Q)$, such that
\[
 \PB_{\sqrt{(1+u)^3}}(V_i),\qquad
 J_0=u\sum_{i=1}^NP_{\C h_i},\qquad
 J_0E=0,\quad\|J_0\|<2u,\quad\ip{e}{\Psi(J_0)e}\ge c.
\]
Each $V_i$ has pairwise distinct nonzero auxiliary modes and rational
real frequencies.
\end{lemma}
\begin{proof}
Set $\phi=(A,f)$, write $r=\sum_j a_j\phi_j$, and let $S$ be the
Gram matrix of $\phi$ in $L^2(\rho)$. Choose, for example,
$z=\frac12\min(c_u,u/(1+\|S\|))$ as in
Lemma~\ref{lem:residual-lift} and define, before
choosing any scalar mass losses,
\[
 \alpha=\frac{\|r\|_\rho}
 {\sqrt{\|r\|_\rho^2+z^{-1}\sum_j|a_j|^2}},\qquad
 N=\lceil\alpha^{-2}\rceil.
\]
These are fixed by $r,z$ and satisfy $0<\alpha<1$.
Take $N$ independent lifts, each using its own scalar coordinates;
choose all their scalar losses to sum to less than $a$.
The product weight $Q$ has mass greater than $1-a$, since
$1-\prod m_j\le\sum(1-m_j)$ for numbers in $[0,1]$.
Normalize each lifted $h$ to unit norm. Its mean is $\alpha e$,
and the centered components are mutually orthogonal.
Its pairing with $E$ vanishes because its mean is proportional to
$r\perp E$. Lemma~\ref{lem:equicorrelated} applies.
The prefix and frequency properties are those of
Lemma~\ref{lem:residual-lift}, unchanged by integrating unused factors.
\end{proof}

\begin{lemma}[A common tolerance for the product step]
\label{lem:product-tolerance}
Let $K\ge1$ be a prefix bound for $A$, and fix $\delta,\tau>0$.
Choose the residual copies of Lemma~\ref{lem:residual-copies} with
mass loss less than $\delta/2$. One can then choose $\eta>0$ so that
directional constructions of tolerance $\eta$ and strength $u$ on
each $(V_i,h_i)$ give, on concatenation,
\[
 \|J-J_0\|<N\eta,\quad \|J\|<3u,\quad
 \PB_{\sqrt{1+14u}}(\hbox{centered new family}),
\]
auxiliary mass loss less than $\delta$, and an angle
$\beta=\|P_F|_E\|<1$ satisfying
\[
 \sqrt{\frac{1+\beta}{1-\beta}}\max(K,K_*)<\max(K,K_*)+\tau,
 \qquad \|P_Fe\|^2>c/2.
\]
Here all base vectors are embedded by the isometry of
Lemma~\ref{lem:centering}.
\end{lemma}
\begin{proof}
First take $Q$ with loss less than $\delta/2$ in
Lemma~\ref{lem:residual-copies}. The base measure is now fixed.
Use Proposition~\ref{prop:directional} with $K_s=1+u$ and
$K_V=\sqrt{(1+u)^3}$ in pairwise disjoint new coordinate blocks.
Since $(1+u)^2-1\le3u$ and $(1+u)^3-1\le4u$,
each centered squared bound is at most
$1+(\sqrt3+2)^2u+\eta$.
Their centered spans are orthogonal by
Lemma~\ref{lem:independent-blocks}, and their information operators
add, so the same bound holds for the concatenation and
$\|J-J_0\|<N\eta$.

Put $k_0=\max(K,K_*)$,
$R=(1+\tau/k_0)^2$, and $b_0=(R-1)/(R+1)\in(0,1)$.
Choose
\[
 N\eta<\min(u,c/2,\delta/2,b_0^2),\qquad
 \eta<[14-(\sqrt3+2)^2]u.
\]
All right sides are positive, since $4\sqrt3<7$.
The norm and centered bounds follow. Total mass loss is at most the
loss of $Q$ plus the sum of the new losses, hence is less than $\delta$.
By Lemma~\ref{lem:information} and $J_0E=0$,
$\beta^2=\|P_E\Psi(J)|_E\|\le N\eta$.
Thus $\beta<b_0$ gives the displayed angle inequality.
The same resolvent estimate gives
$\|P_Fe\|^2=\ip{e}{\Psi(J)e}\ge c-N\eta>c/2$.
\end{proof}

\begin{proposition}[A finite product extension]\label{prop:product-extension}
If $\PB_K(A)$ holds under $\rho$, then for every $\delta,\tau>0$
there are a character polynomial $0<W\le1$ on $\R\times\T^D$
and an ordered list
$B=(e_{b_i}(x)\chi_{k_i}(y))_i$, with rational $b_i$ and pairwise
distinct nonzero $k_i\in\Z^D$, such that
\[
 \int W(x,y)\,dy=m>1-\delta,\quad
 d\sigma=m^{-1}W\,d\rho\,dy,
\]
\[
 \begin{gathered}
 \PB_{K'}(A,B)\text{ in }L^2(\sigma)
 \quad\hbox{for some }K'<\max(K,K_*)+\tau,\\
 d_\sigma(f,(A,B))^2<(1-c/2)d_\rho(f,A)^2.
 \end{gathered}
\]
\end{proposition}
\begin{proof}
Apply Lemma~\ref{lem:product-tolerance} and put $W=Q\prod_iW_i$.
Its auxiliary integral is
$m=(\int Q)\prod_iM_i$, independent of $x$.
It is positive and polynomial, and is at most one.
For the new full Gram matrix,
Lemma~\ref{lem:information} gives $C\le G\le(1+3u)C$.
Thus the new family has a squared prefix bound at most
$(1+\|J\|)(1+14u)<(1+3u)(1+14u)<(1+10u)^2$;
the last inequality follows by subtracting
to obtain $3u+58u^2>0$.
Use the bound $K_*$ and Lemma~\ref{lem:angle} with the angle in
Lemma~\ref{lem:product-tolerance}.

Every old inner product is preserved by the constant marginals.
The candidate $P_Ef+P_Fr$ lies in $E+F$, and therefore
\[
 d_\sigma(f,(A,B))^2
 \le\|r-P_Fr\|^2
 =\|r\|_\rho^2(1-\|P_Fe\|^2)
 <(1-c/2)\|r\|_\rho^2.
\]
Each residual family has distinct old auxiliary modes. The frequency
assertions of Proposition~\ref{prop:directional} apply, and fresh
nonzero scalar modes in different applications separate their output
families. Real frequencies remain rational. The order is the specified
concatenation order, not numerical sorting of frequencies.
\end{proof}

\section{Graph averaging and one extension on the line}
\begin{lemma}[Avoiding finitely many integer hyperplanes]\label{lem:integer-direction}
For a finite $F\subset\Z^D\setminus\{0\}$ there is
$\omega\in\Z^D$ with $k\cdot\omega\ne0$ for every $k\in F$.
\end{lemma}
\begin{proof}
For $D=0$, $F$ is empty. Otherwise the polynomials
$p_k(X)=\sum_{i=0}^{D-1}k_iX^i$ are nonzero. A nonzero polynomial has
finitely many roots, by induction on degree and division by a linear
factor at each root. Choose an integer $b$ outside the union of their
finite root sets and put $\omega=(1,b,\ldots,b^{D-1})$.
\end{proof}

\begin{lemma}[A mean-square vanishing statement]\label{lem:mean-square}
Let $\mu$ be finite and singleton-null on $\R$, $g\in L^1(\mu)$,
and let $l\in\Z\setminus\{0\}$. Then
\[
 \frac1M\sum_{j=1}^M
 \left|\int g(x)e_{jl}(x)\,d\mu(x)\right|^2\longrightarrow0.
\]
\end{lemma}
\begin{proof}
Set $D_M(t)=M^{-1}\sum_{j=1}^M e^{2\pi ijt}$.
The geometric sum gives
$|D_M(t)|\le2/(M|1-e^{2\pi it}|)$ for $t\notin\Z$,
so $D_M(t)\to0$ there. Also $|D_M(t)|\le1$ everywhere.
Expanding the squares and using Fubini
writes the asserted average as
\[
 \iint g(x)\overline{g(y)}D_M(l(x-y))\,d\mu(x)\,d\mu(y).
\]
The majorant $|g(x)g(y)|$ is integrable, with integral $\|g\|_1^2$.
For every $x$ the exceptional $y$ lie in the countable set
$x+l^{-1}\Z$, which is null. This relation is measurable, being a
countable union of closed lines. Fubini makes it product-null.
Dominated convergence proves the limit.
\end{proof}

\begin{lemma}[One subsequence for finitely many sequences]\label{lem:common-subsequence}
If finitely many complex sequences $a_j^{(r)}$ satisfy
$M^{-1}\sum_{j=1}^M|a_j^{(r)}|^2\to0$, there are strictly increasing
positive integers $j_n$ with $a_{j_n}^{(r)}\to0$ for every $r$.
\end{lemma}
\begin{proof}
The nonnegative sequence $A_j=\sum_r|a_j^{(r)}|^2$ has Ces\`aro
averages tending to zero. For every $N$ and $\epsilon>0$ some $j>N$
has $A_j<\epsilon$; otherwise its average for $M>N$ is at least
$\epsilon(M-N)/M$, a contradiction.
Recursively choose $j_n>j_{n-1}$ with $A_{j_n}<1/(n+1)$.
Each squared modulus is bounded by this sum. The empty family is
handled by $j_n=n+1$.
\end{proof}

\begin{lemma}[Simultaneous graph averaging]\label{lem:graph}
Let $\mu$ be finite and singleton-null, and let finitely many functions
on $\R\times\T^D$ have the form
$G(x,y)=\sum_k g_k(x)\chi_k(y)$ with finite supports and
$g_k\in L^1(\mu)$. For any additional finite
$L\subset\Z^D\setminus\{0\}$ there exist $\omega\in\Z^D$ and
$j_n\uparrow\infty$ such that
\[
 k\cdot\omega\ne0\ (k\in L),\qquad
 \int G(x,j_n\omega x)\,d\mu\longrightarrow
 \iint G(x,y)\,dy\,d\mu
\]
for all of the functions simultaneously.
\end{lemma}
\begin{proof}
Apply Lemma~\ref{lem:integer-direction} to $L$ together with every
nonzero mode in every support. A term with mode $k\ne0$ restricts
to $g_k(x)e_{j(k\cdot\omega)}(x)$.
Lemmas~\ref{lem:mean-square} and \ref{lem:common-subsequence} make
all these integrals tend to zero along one subsequence.
Zero-mode integrals are unchanged and equal the torus integrals by
Lemma~\ref{lem:character-algebra}. If $D=0$ only zero modes occur.
\end{proof}

\begin{lemma}[Eventual frequency separation]\label{lem:separate-frequencies}
For a finite ordered list of rational $b_i$ and pairwise distinct
nonzero integers $l_i$, the frequencies $b_i+jl_i$, in that same
order, are distinct and avoid any given finite set of real old
frequencies for all sufficiently large integers $j$.
\end{lemma}
\begin{proof}
Each collision with an old frequency solves a nonconstant affine
equation in $j$, and each collision between two new frequencies has
nonzero slope $l_i-l_{i'}$. There is at most one real solution per
equation. Finitely many equations have only finitely many exceptional
integer solutions. All the new frequencies are rational.
\end{proof}

\begin{proposition}[One extension on the line]\label{prop:line-extension}
Let $\mu$ be a singleton-null probability, let $0<v\le1$ be measurable,
and put $\nu=v\mu$. Let $A$ be a nonempty ordered list of distinct
rational exponentials with $\PB_K(A)$ under $\nu$, and let $f$ be a
rational exponential not in that list. For $\delta,\tau>0$ there is
a positive character polynomial $V\le1$ on $\R$ and an ordered
extension $A'$ by distinct rational exponentials such that
\[
 \int(1-V)\,d\mu<\delta,\qquad
 \PB_{\max(K,K_*)+\tau}(A')\text{ under }V\nu,
\]
\[
 d_{V\nu}(f,A')^2\le(1-c/4)d_\nu(f,A)^2.
\]
\end{proposition}
\begin{proof}
Lemma~\ref{lem:weighted-measures} gives $0<\nu(\R)\le1$ and
singleton-nullity. Normalize to $\rho=\nu/\nu(\R)$.
Apply Proposition~\ref{prop:product-extension} with tolerances
$\delta/2,\tau/2$, obtaining $W,B,m,\sigma$ and a strict prefix
bound $K'<\max(K,K_*)+\tau/2$.
Use Lemma~\ref{lem:graph} with
\[
 G\in\{1-W\}\cup
 \{vW\overline p q:p,q\in(A,B,f)\},\quad
 L=\{k_i\}_i\cup\{k_i-k_{i'}:i\ne i'\}.
\]
All coefficients are integrable: the character polynomials are
bounded and $0<v\le1$ on a probability space.
Set $V_n(x)=W(x,j_n\omega x)$ and replace each member
$e_{b_i}\chi_{k_i}$ of $B$ by $e_{b_i+j_n k_i\cdot\omega}$.
Call the appended ordered list $A'_n$.
Lemma~\ref{lem:separate-frequencies} applies because of the chosen
nonzero, distinct slopes. The functions $V_n$ are positive rational
character polynomials and are at most one.

The product family $(A,B,f)$ has positive definite Gram matrix.
Indeed a weighted zero combination is also zero in the reference
product measure because $W>0$. Haar integration separates its
distinct nonzero auxiliary modes, forcing every $B$ coefficient to
vanish since its real exponential has modulus one. The zero-mode
combination of $(A,f)$ then vanishes, and
Lemma~\ref{lem:real-independence} gives its independence.
Let $s=m\nu(\R)>0$. The Gram matrices on the graph converge to
$s$ times the corresponding $\sigma$ Gram matrices, entry by entry,
and $\int(1-V_n)\,d\mu\to1-m<\delta/2$.
Lemma~\ref{lem:gram-stability}, with sufficiently small relative error,
gives the bound $\max(K,K_*)+\tau$ eventually, using the strict slack
in $K'$ and invariance of prefix bounds under scaling.
The distance formula gives
\[
 d_{V_n\nu}(f,A'_n)^2
 \longrightarrow s\,d_\sigma(f,(A,B))^2
 <(1-c/2)d_\nu(f,A)^2,
\]
because $m\le1$ and distances under $\rho$ scale by $\nu(\R)$.
The old distance is positive by independence. Hence the weaker
factor $1-c/4$ holds eventually. Take one index after all the finitely
many mass, prefix, distance and frequency thresholds.
\end{proof}

\section{Ordered stages and the limiting weighted basis}
\begin{lemma}[A recurrent rational schedule]\label{lem:schedule}
There is a sequence $(q_n)$ of rationals such that each rational
occurs infinitely often. For each rational $q$, its number of
occurrences among $q_0,\ldots,q_{n-1}$ tends to infinity.
\end{lemma}
\begin{proof}
Fix an enumeration $r_0,r_1,\ldots$ of $\Q$ and concatenate the finite
rounds $(r_0)$, $(r_0,r_1)$, $(r_0,r_1,r_2)$, and so on.
Round $m$ starts at $T_m=m(m+1)/2$ and ends before $T_{m+1}$.
For each $k$, the positions $T_m+k$, $m\ge k$, contain $r_k$,
are distinct and tend to infinity. Any prescribed finite number of
these positions is eventually included in the counted initial segment.
\end{proof}

\begin{lemma}[Positive products and convergence of weights]
\label{lem:positive-product}
Let $\mu$ be a probability and $V_n$ measurable with $0<V_n\le1$
pointwise and $\sum_n\int(1-V_n)\,d\mu<\infty$.
Set $v_0=1$ and $v_n=\prod_{j<n}V_j$. Then $v_n\downarrow w$,
where $w>0$ almost everywhere, and $v_n\to w$ in $L^1(\mu)$.
One may modify $w$ on a null set to have $0<w\le1$ everywhere.
\end{lemma}
\begin{proof}
Monotonicity gives a pointwise limit in $[0,1]$, which is measurable.
Tonelli gives $\sum_n(1-V_n(x))<\infty$ almost everywhere.
At such a point $V_n(x)\to1$ and eventually $V_n(x)\ge1/2$.
Since $-\log t=\int_t^1s^{-1}\,ds\le2(1-t)$ on $[1/2,1]$,
the sum of $-\log V_n(x)$ is finite; the finitely many initial
terms are finite because every factor is positive.
Exponentiating the convergent sum proves $w(x)>0$.
Dominated convergence, with bound one for $|v_n-w|$, proves $L^1$
convergence. Assign value one on the measurable null set $\{w=0\}$.
All comparisons with $v_n$ then remain valid almost everywhere.
\end{proof}

\begin{definition}[Extension states]\label{def:states}
A state consists of a measurable pointwise positive weight $v\le1$
on $\R$, a nonempty finite ordered list of distinct rational
frequencies $a=(a_0,\ldots,a_{d-1})$, and a specified prefix bound for
the associated exponential list $A=(e_{a_i})$ under $v\mu$.
A successor retains the frequency list as an initial
segment and multiplies $v$ by a pointwise positive polynomial $V\le1$.
The stage index and the list length are distinct pieces of data.
\end{definition}

\begin{lemma}[Existence of compatible states]\label{lem:states}
Given $\varepsilon>0$, put
\[
 u=\min(1,\varepsilon)/100,\quad
 \delta_n=2^{-n-1},\quad
 \tau_n=(\varepsilon-10u)2^{-n-2},\quad
 K_n=K_*+\sum_{j<n}\tau_j.
\]
For a singleton-null probability $\mu$, there are states $(v_n,A_n)$
with $v_0=1$, $A_0=(e_0)$, and bound $K_n$.
Writing $a_n$ for their frequency lists, $a_{n+1}$ extends $a_n$
as an initial segment, and $v_{n+1}=v_nV_n$, where
$\int(1-V_n)\,d\mu<\delta_n$.
With $\vartheta=1-c/4\in(0,1)$, the scheduled target satisfies
\[
 d_{v_{n+1}\mu}(e_{q_n},A_{n+1})^2
 \le\vartheta\,d_{v_n\mu}(e_{q_n},A_n)^2.
\]
\end{lemma}
\begin{proof}
$\varepsilon-10u>0$, so all tolerances are positive.
The initial one-vector family has bound one, hence bound $K_0$.
If $q_n$ is already in the frequency list, set $V_n=1$ and keep the
state; its target distance is zero, and its old bound is at most
$K_{n+1}$. Otherwise apply Proposition~\ref{prop:line-extension}.
The hypothesis that the target is outside the old span follows from
Lemma~\ref{lem:real-independence}, because $v_n\mu$ is nonzero and
singleton-null. As $K_n\ge K_*$, the new bound is
$K_n+\tau_n=K_{n+1}$. The new polynomial is pointwise positive and
at most one, preserving the state conditions.
Induction with a choice of a witness at each step constructs the states.
No decision about arbitrary closed subspace membership is needed.
\end{proof}

\begin{lemma}[Density of the union in the limiting measure]
\label{lem:limit-density}
For the states in Lemma~\ref{lem:states}, let $w$ be the weight from
Lemma~\ref{lem:positive-product}. Then the union of their spans is
dense in $L^2(w\mu)$, and their lengths $d_n$ are unbounded.
\end{lemma}
\begin{proof}
The multiplier losses are summable. For a fixed rational $q$, let
$N_n(q)$ count visits to $q$ in the first $n$ stages.
Lemma~\ref{lem:distance-monotone} gives nonincrease at every other
stage and after passage from $v_n$ to $w$. The visited stage either
has zero distance or supplies the contraction in Lemma~\ref{lem:states}.
Induction therefore yields
\[
 d_{w\mu}(e_q,A_n)^2
 \le d_{v_n\mu}(e_q,A_n)^2
 \le\vartheta^{N_n(q)}d_\mu(e_q,A_0)^2\longrightarrow0.
\]
Use Lemma~\ref{lem:schedule} and $0<\vartheta<1$.
Each rational exponential lies in the closure of the union, so
Lemma~\ref{lem:real-totality} proves density.
If the nondecreasing integer lengths were bounded, they would
eventually be constant. Compatibility would then make all subsequent
lists identical. Their finite-dimensional span is closed, so density
would make it all of $L^2(w\mu)$.
But $w\mu$ is nonzero and singleton-null, and
Lemma~\ref{lem:real-independence} provides independent families of
arbitrarily large size in that space, a contradiction.
\end{proof}

\begin{lemma}[Extracting the ordered limit list]\label{lem:ordered-limit}
If finite lists $A_n$ extend one another and their lengths $d_n$
are unbounded, there is a unique sequence whose first $d_n$ entries
are $A_n$ for every $n$. If each list has distinct entries, this
sequence is injective.
\end{lemma}
\begin{proof}
For entry $k$, choose any $n$ with $d_n>k$ and use its $k$th entry.
Two choices agree in any later list containing both, so the definition
is independent of the choice and has the required prefixes.
These prefixes exhaust all entries, proving uniqueness.
For two distinct indices choose one list long enough to contain both;
its injectivity distinguishes their entries.
\end{proof}

\begin{lemma}[Passing prefix inequalities to the limit]
\label{lem:limit-prefix}
For the limiting sequence $(e_{t_k})$ extracted from the states,
every finite initial family in $L^2(w\mu)$ has prefix bound
\[
 K_\infty=K_*+\sum_{n=0}^\infty\tau_n
 =1+10u+\frac{\varepsilon-10u}{2}<1+\varepsilon.
\]
\end{lemma}
\begin{proof}
Fix a finite coefficient vector and its cut, giving bounded
character polynomials $p$ and $p_{\rm cut}$.
For all sufficiently large stages both are represented in $A_n$
with exactly those coefficients, padded by zeros.
The state inequality gives
$\|p_{\rm cut}\|_{v_n\mu}^2\le K_\infty^2\|p\|_{v_n\mu}^2$.
By Lemmas~\ref{lem:positive-product} and \ref{lem:weighted-measures},
both integrals converge to their values under $w\mu$. Taking limits
proves the desired finite inequality.
The geometric sum of the $\tau_n$ is
$(\varepsilon-10u)/2$, and $10u<\varepsilon$ gives the strict bound.
\end{proof}

\begin{theorem}[Weighted exponential basis]\label{thm:weighted}
For a singleton-null Borel probability $\mu$ on $\R$ and
$\varepsilon>0$, there are a measurable weight $0<w\le1$ and an
injective sequence of rational times $t_k$ for which $(e_{t_k})$
is a Schauder basis of $L^2(w\mu)$ with a uniform prefix bound
less than $1+\varepsilon$.
\end{theorem}
\begin{proof}
Take the states of Lemma~\ref{lem:states}. The limiting weight and
density are supplied by Lemmas~\ref{lem:positive-product} and
\ref{lem:limit-density}; the latter also proves unbounded lengths.
Apply Lemma~\ref{lem:ordered-limit} to the frequency lists to obtain
the injective ordered sequence,
and Lemma~\ref{lem:real-independence} gives independence of all its
finite initial families under $w\mu$.
Use Lemma~\ref{lem:limit-prefix} and then
Lemma~\ref{lem:basis-criterion}. The optional null-set modification
of $w$ has no effect on the space or any bound.
\end{proof}

\clearpage
\section{A scalar measure for the given unitary group}
\subsection{The orbit coefficient and its positive approximants}
\begin{lemma}[Group identities and the orbit coefficient]\label{lem:group}
The composition law and invertibility of the $U_t$ imply
$U_0=I$ and $U_{-t}=U_t^{-1}$.
For a unit vector $v$, let $\phi(t)=\ip{v}{U_tv}$.
Then $\phi$ is continuous, $\phi(0)=1$, $|\phi|\le1$, and
\[
 \phi(-t)=\overline{\phi(t)},\qquad
 \ip{U_sv}{U_tv}=\phi(t-s).
\]
\end{lemma}
\begin{proof}
Cancel $U_0$ in $U_0U_0=U_0$ and then use the composition law at
$t+(-t)=0$. Strong continuity and continuity of the inner product
give continuity of $\phi$. Unitarity and Cauchy--Schwarz give its
bound. Moving the unitary $U_s$ from both arguments gives the last
identity; conjugate symmetry gives the other one.
\end{proof}

\begin{lemma}[The triangular-window integral]\label{lem:triangle-integral}
For $T>0$ and a function $g:\R\to\C$ continuous on $[-T,T]$,
\[
 \int_0^T\int_0^T g(t-s)\,dt\,ds
 =\int_{-T}^T(T-|r|)g(r)\,dr.
\]
\end{lemma}
\begin{proof}
Translate the inner variable to $r=t-s$. The double integral becomes
the integral of $g(r)$ over $0\le s\le T$, $-s\le r\le T-s$.
Fubini is valid since this is a bounded integrand on a bounded set.
For fixed $r\in[-T,T]$, the admissible $s$ form
$[0,T]\cap[-r,T-r]$, of length $T-r$ if $r\ge0$ and $T+r$ if $r\le0$.
Outside that interval there are no admissible $s$. This proves the formula.
\end{proof}

\begin{lemma}[Positivity from a vector integral]\label{lem:positive-density}
For $T>0$, put $f_T(t)=(1-|t|/T)_+\phi(t)$ and
$q_T(x)=\int_\R f_T(t)e^{-2\pi itx}\,dt$.
Then $f_T$ is continuous and integrable, and
\[
 q_T(x)=\frac1T
 \left\|\int_0^T e^{-2\pi ixs}U_sv\,ds\right\|^2\ge0.
\]
In particular $q_T$ is real-valued and measurable.
\end{lemma}
\begin{proof}
$f_T$ is continuous, bounded by one, and supported in $[-T,T]$.
The vector integrand on the right is continuous with norm one, so
its Bochner integral on $[0,T]$ exists.
Integrating the inner product in both variables gives its squared norm:
integration against a fixed vector is a continuous linear functional
in the second variable, and conjugate symmetry handles the first.
The double integrand is
$e^{-2\pi ix(t-s)}\ip{U_sv}{U_tv}$.
Apply Lemmas~\ref{lem:group} and \ref{lem:triangle-integral} and divide
by $T$ to obtain $q_T$. Measurability, indeed continuity, of $q_T$
also follows directly by dominated convergence with majorant $|f_T|$.
\end{proof}

\begin{lemma}[Mass one before using inversion]\label{lem:density-mass}
For every $T>0$, $q_T\in L^1(\R)$ and $\int q_T=1$.
The probability $d\mu_T=q_T(x)\,dx$ satisfies
$\int e_t\,d\mu_T=f_T(t)$ for every real $t$.
\end{lemma}
\begin{proof}
For $a>0$, the Gaussian Fourier formula and Fubini give
\[
 \int_\R q_T(x)e^{-\pi a x^2}\,dx
 =\int_\R f_T(t)a^{-1/2}e^{-\pi t^2/a}\,dt.
\]
For Fubini, the absolute double integral is bounded by
$\|f_T\|_1\int e^{-\pi a x^2}\,dx<\infty$.
After substituting $t=\sqrt a\,s$, the right side tends to
$f_T(0)\int e^{-\pi s^2}\,ds=1$ by dominated convergence.
On the left the integrands are nonnegative and increase to $q_T$
along any sequence $a\downarrow0$. Monotone convergence gives
$\int q_T=1$, and its reality and nonnegativity give $L^1$ integrability.
Now both $f_T$ and its Fourier transform $q_T$ are integrable.
Mathlib's Fourier inversion theorem, applied to continuous $f_T$,
gives the stated character identity. No inversion was used to
deduce integrability of $q_T$.
\end{proof}

\begin{lemma}[A limiting probability from character limits]
\label{lem:character-limit}
If probabilities $\mu_n$ on $\R$ satisfy
$\int e_t\,d\mu_n\to\phi(t)$ for every $t$, and $\phi$ is continuous
at zero, there is a Borel probability $\mu$ with
$\int e_t\,d\mu=\phi(t)$ for every $t$.
\end{lemma}
\begin{proof}
In Mathlib's convention their characteristic functions converge to
$s\mapsto\phi(s/(2\pi))$, continuous at zero.
The Mathlib tightness theorem for such characteristic-function
limits applies. Prokhorov compactness gives a cluster point $\mu$
of the sequence in the weak topology on probabilities.
For every fixed $t$, integration of the bounded continuous function
$e_t$ is continuous in that topology. Its value at $\mu$ is a cluster
point of the numerical sequence of integrals, and must equal its
limit $\phi(t)$. Indeed two distinct values in the Hausdorff space
$\C$ have disjoint neighborhoods, incompatible with both convergence
and clustering. No limiting probability is assumed in advance.
\end{proof}

\begin{proposition}[Orbit spectral measure]\label{prop:orbit-measure}
For a unit vector $v$ in the given group there is a Borel probability
$\mu$ on $\R$ such that
$\ip{U_sv}{U_tv}=\int e_{t-s}\,d\mu$ for all $s,t$.
\end{proposition}
\begin{proof}
Take the probabilities $\mu_{n+1}$ from Lemma~\ref{lem:density-mass}.
For every fixed $t$, $f_{n+1}(t)\to\phi(t)$.
Apply Lemma~\ref{lem:character-limit} and then Lemma~\ref{lem:group}.
Using $n+1$ ensures every window parameter is positive.
\end{proof}

\subsection{The cyclic isometry and weighted transport}
\begin{lemma}[Equal Gram data on dense spans]\label{lem:equal-gram}
Suppose families $(f_i)_{i\in J}$ and $(g_i)_{i\in J}$ have dense
spans in Hilbert spaces $\mathcal H_1,\mathcal H_2$, and
$\ip{f_i}{f_j}=\ip{g_i}{g_j}$ for all $i,j$.
There is a unique linear isometry equivalence sending every $f_i$
to $g_i$. Independence of either family is not required.
\end{lemma}
\begin{proof}
For every finitely supported coefficient vector the two synthesized
squared norms are equal, by the double-sum Gram formula.
Thus a zero combination on one side is zero on the other, and the
map between algebraic spans is well defined and isometric.
For $x_n$ in the first span with $x_n\to x$, define its image as the
limit of the images of $x_n$. They are Cauchy; equality of norms
makes the limit independent of approximants and preserves linearity
and isometry. Its range is closed: a convergent sequence in the
range lifts to a Cauchy sequence in the complete domain.
The range contains the dense second span, hence is all of $\mathcal H_2$.
Density also proves uniqueness.
\end{proof}

\begin{lemma}[Cyclic realization and intertwining]\label{lem:cyclic-realization}
If $v$ is a cyclic unit vector and $\mu$ is as in
Proposition~\ref{prop:orbit-measure}, there is a linear isometry
equivalence $S:L^2(\mu)\to H$ satisfying
$Se_t=U_tv$ and $S(e_tF)=U_tSF$ for every $F$ and $t$.
\end{lemma}
\begin{proof}
The real exponentials have dense span by
Lemma~\ref{lem:real-totality}, and cyclicity gives the other dense
span. Their Gram matrices agree by Proposition~\ref{prop:orbit-measure},
so Lemma~\ref{lem:equal-gram} applies.
Multiplication by $e_t$ is an isometry of $L^2(\mu)$ with inverse
multiplication by $e_{-t}$. On each $e_s$, the desired identity is
$Se_{s+t}=U_tU_sv$. Linearity proves it on their span; boundedness
of both sides extends it by density.
\end{proof}

\begin{lemma}[No point masses]\label{lem:no-point-masses}
If the group has no common eigenvector, the measure in
Lemma~\ref{lem:cyclic-realization} is singleton-null.
\end{lemma}
\begin{proof}
If $\mu(\{x\})>0$, then $h=\mathbf1_{\{x\}}\in L^2(\mu)$ has positive
squared norm $\mu(\{x\})$. For all $t$,
$e_th=e_t(x)h$, so $Sh\ne0$ is a common eigenvector by intertwining.
This contradicts the hypothesis. Every singleton therefore has measure zero.
\end{proof}

\begin{lemma}[The positive-weight isometry]\label{lem:weight-isometry}
For a finite $\mu$ and a measurable $0<w\le1$ almost everywhere,
multiplication by $\sqrt w$ is a complex-linear isometry equivalence
$L^2(w\mu)\to L^2(\mu)$.
\end{lemma}
\begin{proof}
Choose the representatives of Lemma~\ref{lem:weighted-measures}.
The weighted norm identity gives
$\int|\sqrt w\,F|^2\,d\mu=\int|F|^2w\,d\mu$.
It proves square integrability, well-definedness on equivalence
classes, and isometry; scalar multiplication and addition are
preserved pointwise.
For $G\in L^2(\mu)$ put $F=G/\sqrt w$ on $\{w>0\}$ and zero
elsewhere. This is measurable and
$\int|F|^2w\,d\mu=\int|G|^2\,d\mu<\infty$.
Thus it is a valid inverse, with both identities holding almost
everywhere. No essential bound on $1/\sqrt w$ is needed or asserted.
\end{proof}

\section{Proof of the target theorem}
\begin{proof}[Proof of Theorem~\ref{thm:main}]
Suppose first that $b_n=U_{t_n}g$ is a Schauder basis.
Its span is contained in the orbit span, so $g$ is cyclic.
All its norms equal $\|g\|$, so Lemma~\ref{lem:weak-null} applies.
For a common eigenvector $a\ne0$, unitarity gives
$U_ta=\lambda_ta$ with $|\lambda_t|=1$.
Consequently $|\ip{a}{U_{t_n}g}|=|\ip{a}{g}|$, which must be zero
by weak nullity. Then $a$ is orthogonal to the dense basis span,
contradicting $a\ne0$.

Conversely, a cyclic vector is nonzero because $H$ is nonzero;
dividing it by its norm preserves its orbit span. Choose such a cyclic
unit vector $v$. Proposition~\ref{prop:orbit-measure} and
Lemma~\ref{lem:cyclic-realization} give $\mu$ and $S$.
Lemma~\ref{lem:no-point-masses} makes $\mu$ singleton-null.
Apply Theorem~\ref{thm:weighted}, for example with $\varepsilon=1$.
Since $\int w\,d\mu\le1$, the function $\sqrt w$ belongs to $L^2(\mu)$.
Compose its positive-weight isometry from
Lemma~\ref{lem:weight-isometry} with $S$.
It sends each basis vector $e_{t_k}$ to
$S(\sqrt w\,e_{t_k})=U_{t_k}(S\sqrt w)$.
Lemma~\ref{lem:basis-transport} gives a Schauder basis in this orbit.
The times are rational and injective, hence distinct as real numbers.
Shifting $k=0,1,\ldots$ to $n=1,2,\ldots$ gives exactly the statement.
\end{proof}

\appendix
\section{Mathlib correspondence and the audit changes}\label{sec:mathlib}
The inspected Mathlib revision is
\href{https://github.com/leanprover-community/mathlib4/tree/1ddfaa9548e05b343be9248ae08b73c3425b446e}{\texttt{1ddfaa9548e05b343be9248ae08b73c3425b446e}},
with Lean \texttt{v4.36.0}. The following source pointers identify the
foundational results and the intended interfaces. They are not a claim
that the new lemmas of this document already have Lean proofs.

\subsection{Foundational declarations}
\begingroup\small\raggedright
\begin{enumerate}
\item \ml{GeneralSchauderBasis}, \ml{SchauderBasis}, and \ml{SchauderBasis.exists_norm_proj_le}: continuous coordinates and bounded projections for an already bundled basis. Lemmas~\ref{lem:coordinate-continuity} and \ref{lem:weak-null} clarify the mathematical interface.
\href{https://github.com/leanprover-community/mathlib4/blob/1ddfaa9548e05b343be9248ae08b73c3425b446e/Mathlib/Analysis/Normed/Module/Bases.lean#L104}{\ml{Mathlib/Analysis/Normed/Module/Bases.lean}}.
\item \ml{SchauderBasis.RankOneDecomposition.basis} and \ml{basis_coe}: the projection conditions verified in Lemma~\ref{lem:basis-criterion} preserve the specified exponential sequence.
\href{https://github.com/leanprover-community/mathlib4/blob/1ddfaa9548e05b343be9248ae08b73c3425b446e/Mathlib/Analysis/Normed/Module/Bases.lean#L464}{\ml{Mathlib/Analysis/Normed/Module/Bases.lean}}.
\item \ml{ContinuousLinearEquiv.ofBijective}: the bounded inverse theorem in Lemma~\ref{lem:coordinate-continuity}; its coefficient-space domain is proved complete in Lemma~\ref{lem:coefficient-space}.
\href{https://github.com/leanprover-community/mathlib4/blob/1ddfaa9548e05b343be9248ae08b73c3425b446e/Mathlib/Analysis/Normed/Operator/Banach.lean#L439}{\ml{Mathlib/Analysis/Normed/Operator/Banach.lean}}.
\item \ml{InnerProductSpace.toDual}: representation of the coordinate functionals in Lemma~\ref{lem:weak-null}. Completeness and the complex Hilbert structure are hypotheses of the application.
\href{https://github.com/leanprover-community/mathlib4/blob/1ddfaa9548e05b343be9248ae08b73c3425b446e/Mathlib/Analysis/InnerProductSpace/Dual.lean#L136}{\ml{Mathlib/Analysis/InnerProductSpace/Dual.lean}}.
\item \ml{Submodule.orthogonalDecomposition}: orthogonal sums use \ml{WithLp 2}, with the Hilbert sum-of-squares metric. Orthogonal projection and the orthogonal-complement density criterion are used throughout.
\href{https://github.com/leanprover-community/mathlib4/blob/1ddfaa9548e05b343be9248ae08b73c3425b446e/Mathlib/Analysis/InnerProductSpace/ProdL2.lean#L88}{\ml{Mathlib/Analysis/InnerProductSpace/ProdL2.lean}}.
\item \ml{Matrix.gram}, \ml{star_dotProduct_gram_mulVec}, and \ml{posDef_gram_iff_linearIndependent}: these agree with our Gram convention. Lemma~\ref{lem:gram-prefix} gives the finite prefix predicate.
\href{https://github.com/leanprover-community/mathlib4/blob/1ddfaa9548e05b343be9248ae08b73c3425b446e/Mathlib/Analysis/InnerProductSpace/GramMatrix.lean#L40}{\ml{Mathlib/Analysis/InnerProductSpace/GramMatrix.lean}}.
\item \ml{Matrix.Norms.L2Operator} and \ml{Matrix.l2_opNorm_def}: the matrix norms in this document. The ordinary entrywise matrix norm would not express the estimates proved here.
\href{https://github.com/leanprover-community/mathlib4/blob/1ddfaa9548e05b343be9248ae08b73c3425b446e/Mathlib/Analysis/CStarAlgebra/Matrix.lean#L195}{\ml{Mathlib/Analysis/CStarAlgebra/Matrix.lean}}.
\item \ml{MatrixOrder}, \ml{Matrix.isStrictlyPositive_iff_posDef}, and the finite matrix functional calculus: Loewner order and square roots are supported. Our inverse variational and averaging arguments are proved directly in Lemma~\ref{lem:inverse-variational}.
\href{https://github.com/leanprover-community/mathlib4/blob/1ddfaa9548e05b343be9248ae08b73c3425b446e/Mathlib/Analysis/Matrix/Order.lean#L54}{\ml{Mathlib/Analysis/Matrix/Order.lean}}.
\item \ml{LinearMap.IsSymmetric.eigenvectorBasis}: finite self-adjoint diagonalization. Every spectral calculation here is finite-dimensional, or explicitly reduced to the finite range of an information operator.
\href{https://github.com/leanprover-community/mathlib4/blob/1ddfaa9548e05b343be9248ae08b73c3425b446e/Mathlib/Analysis/InnerProductSpace/Spectrum.lean#L296}{\ml{Mathlib/Analysis/InnerProductSpace/Spectrum.lean}}.
\item \ml{Matrix.PosDef.fromBlocks}$_{11}$ and \ml{Matrix.PosDef.fromBlocks}$_{22}$: finite Schur complements with a positive invertible pivot. The corresponding square completions and specialized identities are included in the proofs.
\href{https://github.com/leanprover-community/mathlib4/blob/1ddfaa9548e05b343be9248ae08b73c3425b446e/Mathlib/LinearAlgebra/Matrix/PosDef.lean#L698}{\ml{Mathlib/LinearAlgebra/Matrix/PosDef.lean}}.
\item \ml{MeasureTheory.tendsto_integral_of_dominated_convergence}: applied with the majorants specified in Lemmas~\ref{lem:density-transform}, \ref{lem:mean-square}, \ref{lem:positive-product}, and \ref{lem:density-mass}.
\href{https://github.com/leanprover-community/mathlib4/blob/1ddfaa9548e05b343be9248ae08b73c3425b446e/Mathlib/MeasureTheory/Integral/DominatedConvergence.lean#L57}{\ml{Mathlib/MeasureTheory/Integral/DominatedConvergence.lean}}.
\item \ml{MeasureTheory.integral_prod}, together with Tonelli and monotone convergence for nonnegative integrands: finite product densities, graph averaging, and Gaussian regularization. All product measures are finite or sigma-finite.
\href{https://github.com/leanprover-community/mathlib4/blob/1ddfaa9548e05b343be9248ae08b73c3425b446e/Mathlib/MeasureTheory/Integral/Prod.lean#L443}{\ml{Mathlib/MeasureTheory/Integral/Prod.lean}}.
\item \ml{MeasureTheory.integral_withDensity_eq_integral_smul} and its variants: the norm and inner-product identities for weighted measures. The definition uses a nonnegative extended-real density.
\href{https://github.com/leanprover-community/mathlib4/blob/1ddfaa9548e05b343be9248ae08b73c3425b446e/Mathlib/MeasureTheory/Integral/Bochner/ContinuousLinearMap.lean#L251}{\ml{Mathlib/MeasureTheory/Integral/Bochner/ContinuousLinearMap.lean}}.
\item \ml{MeasureTheory.memLp_two_iff_integrable_sq_norm}, the \ml{L2} inner-product formula, and completeness of \ml{Lp}: these support Lemma~\ref{lem:weight-isometry}. Its inverse is justified in the weighted squared norm.
\href{https://github.com/leanprover-community/mathlib4/blob/1ddfaa9548e05b343be9248ae08b73c3425b446e/Mathlib/MeasureTheory/Function/L2Space.lean#L45}{\ml{Mathlib/MeasureTheory/Function/L2Space.lean}}.
\item \ml{integral_inner}: integration of Hilbert-valued orbit vectors in Lemma~\ref{lem:positive-density}; the continuous integrands on compact intervals are Bochner integrable.
\href{https://github.com/leanprover-community/mathlib4/blob/1ddfaa9548e05b343be9248ae08b73c3425b446e/Mathlib/MeasureTheory/Function/L2Space.lean#L89}{\ml{Mathlib/MeasureTheory/Function/L2Space.lean}}.
\item \ml{MeasureTheory.NullSingletonClass} and \ml{Set.Countable.measure_zero}: precisely the no-point-mass assumption used in Lemma~\ref{lem:mean-square}.
\href{https://github.com/leanprover-community/mathlib4/blob/1ddfaa9548e05b343be9248ae08b73c3425b446e/Mathlib/MeasureTheory/Measure/Typeclasses/NullSingletonClass.lean#L28}{\ml{Mathlib/MeasureTheory/Measure/Typeclasses/NullSingletonClass.lean}}.
\item \ml{MeasureTheory.exists_accPt_of_nullSingletonClass}: a positive-measure subset of the real line has an accumulation point. The required separability of the subset is automatic here.
\href{https://github.com/leanprover-community/mathlib4/blob/1ddfaa9548e05b343be9248ae08b73c3425b446e/Mathlib/MeasureTheory/Measure/Typeclasses/NullSingletonClass.lean#L88}{\ml{Mathlib/MeasureTheory/Measure/Typeclasses/NullSingletonClass.lean}}.
\item \ml{linearIndependent_monoidHom}: independence of distinct characters on the additive real group, encoded by its multiplicative type synonym, in Lemma~\ref{lem:pointwise-independence}.
\href{https://github.com/leanprover-community/mathlib4/blob/1ddfaa9548e05b343be9248ae08b73c3425b446e/Mathlib/LinearAlgebra/LinearIndependent/Basic.lean#L516}{\ml{Mathlib/LinearAlgebra/LinearIndependent/Basic.lean}}.
\item \ml{AnalyticOnNhd.eqOn_zero_of_preconnected_of_frequently_eq_zero}: the real-analytic, complex-valued identity principle in Lemma~\ref{lem:real-independence}.
\href{https://github.com/leanprover-community/mathlib4/blob/1ddfaa9548e05b343be9248ae08b73c3425b446e/Mathlib/Analysis/Analytic/IsolatedZeros.lean#L210}{\ml{Mathlib/Analysis/Analytic/IsolatedZeros.lean}}.
\item \ml{AddCircle.span_fourierLp_closure_eq_top}: period one, exponent two, normalized Haar measure. Only the one-variable exhaustion in Lemma~\ref{lem:rare-data} uses this totality result.
\href{https://github.com/leanprover-community/mathlib4/blob/1ddfaa9548e05b343be9248ae08b73c3425b446e/Mathlib/Analysis/Fourier/AddCircle.lean#L259}{\ml{Mathlib/Analysis/Fourier/AddCircle.lean}}.
\item \ml{UnitAddTorus.mFourier}: the finite auxiliary characters. Their moments and the finite Fejer kernel required here are proved in Lemmas~\ref{lem:character-algebra} and \ref{lem:fejer}.
\href{https://github.com/leanprover-community/mathlib4/blob/1ddfaa9548e05b343be9248ae08b73c3425b446e/Mathlib/Analysis/Fourier/AddCircleMulti.lean#L53}{\ml{Mathlib/Analysis/Fourier/AddCircleMulti.lean}}.
\item \ml{MeasureTheory.Measure.ext_of_charFun}: finite positive Borel measures on the complete second-countable real line. Lemma~\ref{lem:density-transform} proves the reduction from a complex density.
\href{https://github.com/leanprover-community/mathlib4/blob/1ddfaa9548e05b343be9248ae08b73c3425b446e/Mathlib/MeasureTheory/Measure/CharacteristicFunction/Basic.lean#L256}{\ml{Mathlib/MeasureTheory/Measure/CharacteristicFunction/Basic.lean}}.
\item \ml{fourier_gaussian_pi}: the positive Gaussian parameter in Lemma~\ref{lem:density-mass} gives exactly the stated Fourier normalization.
\href{https://github.com/leanprover-community/mathlib4/blob/1ddfaa9548e05b343be9248ae08b73c3425b446e/Mathlib/Analysis/SpecialFunctions/Gaussian/FourierTransform.lean#L227}{\ml{Mathlib/Analysis/SpecialFunctions/Gaussian/FourierTransform.lean}}.
\item \ml{Continuous.fourierInv_fourier_eq}: both integrability hypotheses are established before inversion in Lemma~\ref{lem:density-mass}.
\href{https://github.com/leanprover-community/mathlib4/blob/1ddfaa9548e05b343be9248ae08b73c3425b446e/Mathlib/Analysis/Fourier/Inversion.lean#L175}{\ml{Mathlib/Analysis/Fourier/Inversion.lean}}.
\item \ml{MeasureTheory.isTightMeasureSet_of_tendsto_charFun}: applied on the real line to the pointwise characteristic-function limit continuous at zero. It supplies tightness, not a pre-assumed limiting measure.
\href{https://github.com/leanprover-community/mathlib4/blob/1ddfaa9548e05b343be9248ae08b73c3425b446e/Mathlib/MeasureTheory/Measure/LevyConvergence.lean#L43}{\ml{Mathlib/MeasureTheory/Measure/LevyConvergence.lean}}.
\item \ml{isCompact_closure_of_isTightMeasureSet}: weak compactness of the closure of the tight family of Borel probabilities in Lemma~\ref{lem:character-limit}.
\href{https://github.com/leanprover-community/mathlib4/blob/1ddfaa9548e05b343be9248ae08b73c3425b446e/Mathlib/MeasureTheory/Measure/Prokhorov.lean#L553}{\ml{Mathlib/MeasureTheory/Measure/Prokhorov.lean}}.
\item \ml{MeasureTheory.ProbabilityMeasure.tendsto_iff_forall_integral_rclike_tendsto}: bounded continuous complex characters are legitimate weak test functions.
\href{https://github.com/leanprover-community/mathlib4/blob/1ddfaa9548e05b343be9248ae08b73c3425b446e/Mathlib/MeasureTheory/Measure/ProbabilityMeasure.lean#L374}{\ml{Mathlib/MeasureTheory/Measure/ProbabilityMeasure.lean}}.
\item \ml{LinearEquiv.extendOfIsometry}: an available implementation of the dense-span extension proved in Lemma~\ref{lem:equal-gram}. Its use does not require independence of the generating families.
\href{https://github.com/leanprover-community/mathlib4/blob/1ddfaa9548e05b343be9248ae08b73c3425b446e/Mathlib/Analysis/Normed/Operator/Extend.lean#L305}{\ml{Mathlib/Analysis/Normed/Operator/Extend.lean}}.
\end{enumerate}
\endgroup

\subsection{Where the audit findings are addressed}
\begin{enumerate}
\item Definitions~\ref{def:characters}, \ref{def:prefix},
\ref{def:scalar}, and \ref{def:states} retain bounded representatives,
ordered frequencies, scalar data, and stage compatibility explicitly.
No optimal finite basis constant is constructed.
\item Lemmas~\ref{lem:inverse-exhaustion}, \ref{lem:row-matrix},
and \ref{lem:two-projections} separate the changing-dimension inverse
limit and the two finite matrix calculations previously compressed
into the row-splitting argument. Lemma~\ref{lem:row-split} works on
finite Hilbert tuples, with an explicitly given inverse operator.
\item Lemmas~\ref{lem:scalar-numerics}, \ref{lem:scalar-inverse-step},
and \ref{lem:scalar-prefix-step} expose every scalar induction
invariant. Lemmas~\ref{lem:fejer}--\ref{lem:tuning} prove only the
finite smoothing facts needed for exact-gain scalar supply.
\item Lemma~\ref{lem:seed-scales} uses an explicit sequence of changing
finite scalar systems and uniform response estimates.
Lemmas~\ref{lem:column-bounds} and \ref{lem:selected-prefix} control
cuts inside columns as well as between columns.
\item Lemmas~\ref{lem:residual-copies} and \ref{lem:product-tolerance}
fix the order of all choices and a single common tolerance.
Lemma~\ref{lem:gram-stability} transfers prefix inequalities directly,
without continuity of an optimized basis constant.
\item Lemmas~\ref{lem:schedule}--\ref{lem:limit-prefix} separate skipped
stages, list lengths, append order, positivity, density, and limiting
finite inequalities. The ordered sequence is not an arbitrary
re-enumeration of a frequency set.
\item Lemmas~\ref{lem:pointwise-independence} and
\ref{lem:real-independence} use character independence and an analytic
accumulation point instead of iterated derivatives and Vandermonde
bookkeeping. Lemmas~\ref{lem:positive-density}--\ref{lem:character-limit}
use a positive orbit integral and the existing tightness theorem.
\item Lemmas~\ref{lem:equal-gram} and \ref{lem:weight-isometry} give the
two isometry constructions separately, with well-definedness and
surjectivity proved. The latter does not assume the reciprocal weight
is essentially bounded.
\end{enumerate}

\subsection{Correspondence with the target and provenance}
The statement with label \ml{thm:main} is unchanged from
\ml{paper.tex}. In \ml{Showcase.lean}, the composition law plus
invertibility gives the identities in Lemma~\ref{lem:group}.
The definitions of cyclic vector and common eigenvector agree with
those in Section~1. Lemma~\ref{lem:coordinate-continuity} reconciles
the conventional expansion-only definition with Mathlib's continuous
coordinates, and Lemma~\ref{lem:basis-criterion} verifies the data for
its existing basis constructor. The only index change is the shift
between natural numbers starting at zero and positive integers.
The inner-product convention changes notation, not the theorem.

The finite scalar and directional method follows the mathematical
argument in the original paper and its cited source \cite{FLT}.
The reference records provenance; every specialized step needed here
has a proof above. No external Lean implementation has been reused.
This document is a mathematical foundation for a future formalization;
it is not a claim that the proof has been checked by the Lean kernel.

\begin{thebibliography}{9}
\raggedright
\bibitem{FLT} D.~Freeman, L.~Liehr and M.~A.~Taylor,
\emph{Structured coordinate systems in $L^p(\R)$},
\href{https://arxiv.org/html/2610.18709v2}{arXiv:2610.18709v2}, 2026.
\end{thebibliography}
\end{document}
